\documentclass[11pt,letterpaper]{article}
\usepackage[T1]{fontenc}
\usepackage{lmodern}
\usepackage[margin=1in]{geometry}
\usepackage{amsmath,amssymb,amsthm,mathtools,mathrsfs}
\usepackage{microtype}
\usepackage{enumitem}
\usepackage[hidelinks]{hyperref}

\newtheorem{theorem}{Theorem}[section]
\newtheorem{lemma}[theorem]{Lemma}
\newtheorem{proposition}[theorem]{Proposition}
\newtheorem{corollary}[theorem]{Corollary}
\newtheorem{claim}[theorem]{Claim}
\newtheorem*{theoremU}{Theorem U}
\theoremstyle{definition}
\newtheorem{definition}[theorem]{Definition}
\theoremstyle{remark}
\newtheorem{remark}[theorem]{Remark}

\newcommand{\R}{\mathbb R}
\newcommand{\Q}{\mathbb Q}
\newcommand{\Z}{\mathbb Z}
\newcommand{\F}{\mathbb F}
\newcommand{\E}{\mathbb E}
\newcommand{\eps}{\varepsilon}
\newcommand{\ind}{\mathbf 1}
\newcommand{\Tr}{\mathcal T}
\DeclareMathOperator{\Var}{Var}
\DeclareMathOperator{\tw}{tw}
\DeclareMathOperator{\pw}{pw}
\DeclareMathOperator{\td}{td}
\DeclareMathOperator{\bw}{bw}
\DeclareMathOperator{\mmw}{mmw}
\DeclareMathOperator{\route}{route}
\DeclareMathOperator{\live}{live}
\setlength{\emergencystretch}{1.5em}
\setlist[enumerate]{topsep=4pt,itemsep=3pt}

\title{Consequences of constant-factor hardness for uniform sparsest cut:\\
treewidth, bicliques, dense subgraphs and small-set expansion}
\author{Yuwen Zhou\\[2pt] \normalsize University of British Columbia Okanagan\\
\normalsize\texttt{zhouyw07@student.ubc.ca}}
\date{November 2026}

\begin{document}
\maketitle

\begin{abstract}
A September 2026 preprint of OpenAI proves that Uniform Sparsest Cut is
NP-hard to approximate within every constant factor. Before its final
replication step, its proof yields a weighted statement: for every fixed
\(B\), satisfiable formulas give a cut of bounded capacity whose two sides
have demand mass bounded away from zero, while unsatisfiable formulas give
capacity at least \(B\mu(S)(1-\mu(S))\) for every vertex set \(S\). We take
this statement as an imported theorem, called Theorem~U, and derive
NP-hardness results from it. Assuming Theorem~U, treewidth, pathwidth,
treedepth and minimum fill-in are NP-hard to approximate within every
constant factor on bipartite graphs, and so are a family of layout, width,
pebbling and balanced-partitioning problems; Maximum Balanced Biclique and
Maximum Edge Biclique are NP-hard to approximate within \(N^{\gamma}\) for a
fixed \(\gamma>0\); Densest \(k\)-Subgraph, Densest At-Least-\(k\)-Subgraph
and Minimum \(k\)-Cut are NP-hard to approximate within a fixed factor
\(1+\eps_0\); and a gap version of small-set expansion with one fixed
conductance threshold holds on fixed-degree simple graphs. For most of these
problems every-constant hardness was previously known only under the
Small-Set Expansion hypothesis or under stronger running-time assumptions.
The proofs are transfers: products and symmetric balancing of the capacity
system, finite-field incidence graphs that turn capacities into simple
graphs without losing small cuts, the incidence construction of Wu,
Austrin, Pitassi and Liu, a port-expander degree reduction, and an
expander-walk product. All results are conditional on the unrefereed
preprint.
\end{abstract}

\section{Introduction}

In Uniform Sparsest Cut one is given a graph with nonnegative edge
capacities \(c\) on \(N\) vertices and seeks a nonempty proper vertex set
\(S\) minimizing \(c(\partial S)/(|S|(N-|S|))\). The preprint
\cite{oai26} proves that approximating this problem within any fixed
factor is NP-hard. The last step of that proof replicates the vertices of a
weighted capacity system and adds heavy coherence edges. Before that step,
the proof provides a statement of independent interest, which we isolate as
Theorem~U in Section~\ref{sec:prelim}: for every fixed \(B\), one can compute
in polynomial time a capacity system \((V,c,\mu)\) of bounded total capacity
such that
\begin{itemize}[topsep=3pt,itemsep=2pt]
\item for satisfiable formulas, some cut has capacity at most \(21\) and
both sides have \(\mu\)-mass at least \(1/20000\);
\item for unsatisfiable formulas, \(c(\partial S)\ge B\mu(S)(1-\mu(S))\) for
every \(S\subseteq V\).
\end{itemize}
The NO inequality is a Poincar\'e-type inequality with an arbitrarily large
constant, valid at every scale of \(\mu(S)\). Gaps of this shape, obtained
from the Small-Set Expansion (SSE) hypothesis
\cite{rs10,rst12}, are the starting point of known reductions for layout,
width, partitioning, biclique and density problems
\cite{apw12,wapl14,yam18,bfk14,man18}. This note shows that Theorem~U,
although weaker than the SSE hypothesis, suffices for a substantial part of
those consequences.

\paragraph{Status of the results.}
Every result in this note is conditional on Theorem~U. Theorem~U is read off
from Theorem~6.1, Proposition~7.1 and Section~8.1 of \cite{oai26}, an
unrefereed preprint. The hardness proof of \cite{oai26} is not covered by
the Lean formalization accompanying that release, which concerns a separate
integrality-gap result \cite{oailean}. We make no claim about the correctness
of \cite{oai26}. Our contribution consists of the transfers from the
sparsest-cut gap of Theorem~U to the problems listed in
Theorem~\ref{th:main}. All reductions are deterministic.

\subsection{Main results}

\begin{theorem}[Main theorem]\label{th:main}
Assume Theorem~U. Each statement below is witnessed by a deterministic
polynomial-time reduction from 3SAT; consequently none of the approximation
tasks below can be performed in polynomial time unless
\(\mathrm P=\mathrm{NP}\).
\begin{enumerate}[label=\textup{(\Alph*)}]
\item \emph{Treewidth family} (Section~\ref{sec:tw}). For every fixed
\(c>1\) it is NP-hard, on simple bipartite graphs \(H\) with a computable
scale \(s\), to distinguish \(\td(H)\le s\) from \(\tw(H)>cs\). Hence
treewidth, pathwidth and treedepth are NP-hard to approximate within every
constant factor on bipartite graphs, and so are:
minimum fill-in (number of added edges);
branch-width and maximum-matching width;
clique-width, NLC-width, rank-width, Boolean-width, mim-width and
tree-independence number on line graphs of bipartite graphs;
cutwidth, minimum linear arrangement, every fixed \(\ell_p\)-norm of the
prefix cut sizes, carving-width, the vertex-sum layout, interval completion
(total and added edges), and the four directed layout objectives of
\cite{wapl14};
one-shot black pebbling, also on single-sink DAGs of maximum indegree two,
and one-shot black-white pebbling;
minimum bisection, balanced edge and vertex separators, min-sum and min-max
balanced \(k\)-partitioning for each fixed \(k\ge2\) (soundness against
parts of up to \(\beta N/k\) vertices, for any fixed \(1\le\beta<k\)), and
cardinality-normalized small-set edge expansion.
\item \emph{Gap small-set expansion} (Section~\ref{sec:sse}). For every fixed
rational \(0<\delta\le1/4\) there are a fixed \(\eta=\eta(\delta)<1/2\) and a
fixed integer \(d\) such that, for each fixed integer \(j\ge1\), it is
NP-hard, given a simple bipartite \(d\)-regular graph on \(N\) vertices, to
distinguish between the case that some set of exactly \((\delta/j)N\)
vertices has conductance at most \(\eta\), and the case that every nonempty
set of at most \((\delta/j)N\) vertices has conductance at least \(1-\eta\).
A window version, with an arbitrarily large fixed ratio between the NO and
YES conductance bounds, is Theorem~\ref{th:window}.
\item \emph{Bicliques} (Section~\ref{sec:biclique}). There is a fixed
\(\gamma>0\) such that Maximum Balanced Biclique and Maximum Edge Biclique
are NP-hard to approximate within \(N^\gamma\) on bipartite graphs with
\(N\) vertices.
\item \emph{Dense subgraphs and cuts} (Section~\ref{sec:dense}). There is a
fixed \(\eps_0>0\) such that Densest \(k\)-Subgraph, Densest
At-Least-\(k\)-Subgraph and Minimum \(k\)-Cut (with \(k\) part of the input)
are NP-hard to approximate within \(1+\eps_0\) on simple unweighted graphs.
For Densest \(k\)-Subgraph this holds on bipartite graphs of fixed degree,
also when the solution must take \(k/2\) vertices from each shore; for
Minimum \(k\)-Cut it holds on connected graphs.
\end{enumerate}
\end{theorem}

Part~(B) has one fixed threshold \(\eta(\delta)\) close to \(1/2\). The SSE
hypothesis asks for hardness at every prescribed small \(\eta\), so part~(B)
does not imply it; see Section~\ref{sec:open}.

\subsection{Related work and priority}

The following summary reflects literature searches through 18 November 2026
(arXiv, ECCC, DBLP, author pages and proceedings). ``Apparently new'' means
that we did not locate the precise statement; it is not a claim of
exhaustive search.

\paragraph{Treewidth family.}
Every-constant hardness of treewidth, pathwidth, cutwidth, minimum linear
arrangement, the eight layout objectives, interval completion and one-shot
pebbling, including the single-sink indegree-two variant, was proved under
the SSE hypothesis by Austrin, Pitassi and Wu \cite{apw12} and in the journal
paper of Wu, Austrin, Pitassi and Liu \cite{wapl14}, which also covers
minimum fill-in (with an extended abstract at IJCAI~2015 \cite{wapl15}). The
SSE-based results for minimum linear arrangement, balanced separator and
minimum bisection are due to Raghavendra, Steurer and Tulsiani \cite{rst12},
and Yamazaki \cite{yam18} transferred SSE-based hardness to carving-,
branch-, clique-, NLC-, rank- and Boolean-width. Without hypotheses beyond
\(\mathrm P\ne\mathrm{NP}\), Bonnet proved that treewidth is NP-hard to
approximate within \(1.00005\) \cite{bon} (first public on 17 June 2024;
STOC~2025; J.~ACM 2026), Bonnet, Neuen and Soko\l{}owski proved \(1.0003\)
for treedepth \cite{bns25}, and Cao and Sandeep ruled out a PTAS for minimum
fill-in \cite{cs17}. Ambühl, Mastrolilli and Svensson \cite{ams11} ruled out
a PTAS for minimum linear arrangement under a randomized subexponential-time
assumption. For mim-width, constant-factor obstructions were known under
\(\mathrm{NP}\ne\mathrm{ZPP}\) \cite{sv16}, and bounded-width gaps were
proved recently \cite{icalp25,icalp26}; we do not claim the first
constant-factor inapproximability for mim-width. Tree-independence number is
known to be hard to approximate within every constant on unrestricted graphs
\cite{dms21,dfgkm22,zuc07}; our statement is new only on line graphs of
bipartite graphs. For balanced partitioning see Andreev and R\"acke
\cite{ar06} and Bansal et al.\ \cite{bfk14}. The joint promise
``small treedepth versus large treewidth'' appears new; the joint theorem of
\cite{wapl14} names pathwidth on the YES side. Part~(A) follows the
incidence construction of \cite{apw12,wapl14}; the new input is the common
hard family of Theorem~\ref{th:family}, built from Theorem~U. The
contribution is the every-constant NP-hardness, conditional only on
Theorem~U.

\paragraph{Small-set expansion.}
The SSE hypothesis was introduced by Raghavendra and Steurer \cite{rs10};
reductions between expansion problems are in \cite{rst12}, and Raghavendra
and Schramm \cite{rsch14} amplify a gap by random walks when the soundness
is sufficiently larger than the square root of the completeness. The
fixed-degree gap of Theorem~\ref{th:main}(B) appears new in the stated
NP-hard form. The window version (Theorem~\ref{th:window}) is partially
known: \cite[Theorem~3.5 and Corollary~3.6]{rst12} give equal-volume YES
partitions and a NO expansion profile under the SSE hypothesis; the NP-hard
counting-measure form with a prescribed exact density appears new. The
degree reduction in Section~\ref{sec:sse} is related to the zig-zag product
\cite{rvw02}.

\paragraph{Bicliques.}
Peeters proved NP-completeness of Maximum Edge Biclique \cite{pee03}. Feige
\cite{fei02} derived polynomial-factor hardness from average-case hardness
of refuting random 3SAT, Feige and Kogan \cite{fk04} obtained superconstant
factors for both biclique objectives under a stronger running-time
assumption, and Khot \cite{kho06} ruled out a PTAS for Maximum Balanced
Biclique assuming \(\mathrm{NP}\not\subseteq\bigcap_{\eps>0}
\mathrm{BPTIME}(2^{n^\eps})\). Ambühl, Mastrolilli and Svensson
\cite[Theorem~1.4]{ams11} gave a polynomial-factor lower bound for Maximum
Edge Biclique under a randomized subexponential-time assumption, and
Bhangale et al.\ \cite{bhangale16} obtained biclique hardness from a
strengthened Unique Games conjecture. Manurangsi \cite{man17icalp,man18}
proved, under the SSE hypothesis, a base gap between near-half and
arbitrarily small bicliques and amplified it to \(n^{1-\eps}\) by
randomized products. Manurangsi's September 2026 paper \cite{man26} still
records the absence of even a fixed \(1+\eps\) NP-hardness factor for
Maximum Balanced Biclique. Deterministic amplification of biclique gaps by
expander walks on both shores goes back to Goerdt and Lanka
\cite[Section~4]{gl04}, building on derandomized graph products
\cite{afwz95}; the two-shore estimate used in Section~\ref{sec:biclique}
also follows from Ta-Shma and Zadicario \cite[Proposition~22]{tz24}. We do
not present the amplification as new. The new part is the base gap derived
from Theorem~U and the resulting deterministic NP-hardness, which appears
new for both objectives.

\paragraph{Dense subgraphs and cuts.}
Khot \cite{kho06} ruled out a PTAS for Densest \(k\)-Subgraph under the
randomized assumption above; Ghoshal and Lee \cite{gl22} treated the
linear-\(k\) regime under the SSE hypothesis; the best algorithm achieves
\(n^{1/4+\eps}\) \cite{bccfv10}; Manurangsi \cite{man17} proved
almost-polynomial hardness under ETH, reformulated under Almost-ETH in an
October 2026 revision. Laekhanukit, Manurangsi and Trabelsi \cite{lmt26}
still record the absence of a fixed-factor NP-hardness result. Our result
for Densest \(k\)-Subgraph appears new as a fixed-factor NP-hardness theorem.
For Densest At-Least-\(k\)-Subgraph, Bergner (2013 thesis, as attributed in
\cite{man18}) and Manurangsi \cite{man18} proved \(2-\eps\) hardness under
the SSE hypothesis; ordinary fixed-factor NP-hardness appears new. For
Minimum \(k\)-Cut, Saran and Vazirani \cite{sv95} stated fixed
\(1+\eps\) hardness, citing a private communication of Papadimitriou, and
Manurangsi \cite{man18} reports that the proof was not published and proves
\(2-\eps\) hardness under the SSE hypothesis. Our Minimum \(k\)-Cut result is
therefore an explicit derivation, conditional on Theorem~U, not a first
result. The transfers in Section~\ref{sec:dense} are fixed-gap versions of
the SSE-based transfers of \cite{man18}.

\subsection{Techniques and organization}

Section~\ref{sec:prelim} states Theorem~U and three tools: tensorization of
the NO inequality under products, a symmetric balancing construction that
turns the single YES cut into a \(K\)-coloring with exactly equal class
masses, and finite-field incidence matrices whose second singular value is
at most the square root of their degree.

Section~\ref{sec:tw} builds a common hard family (Theorem~\ref{th:family}):
an iterated product gives a \(K\)-ary hierarchy of cheap cuts in YES;
replication with correction edges gives exact vertex counts; incidence
matrices convert capacities into a simple graph while keeping the NO lower
bound for cuts of every size. The width, layout, fill-in and pebbling
consequences then use the incidence construction of \cite{apw12,wapl14}.

Section~\ref{sec:sse} normalizes the balanced system into a positive
semidefinite reversible chain and reduces its degree by port expanders and
incidence matrices; the positive semidefinite inequality
\(\alpha^\top P\beta\le f^\top Pf\) controls sets that meet the two shores
unequally. Section~\ref{sec:biclique} complements a conflict graph to
obtain a biclique gap and amplifies it by an expander-walk product.
Section~\ref{sec:dense} contains elementary transfers from conductance to
dense subgraphs and cuts. Section~\ref{sec:open} lists open problems.

Sections~\ref{sec:tw}, \ref{sec:sse} and~\ref{sec:biclique} depend only on
Section~\ref{sec:prelim}; Section~\ref{sec:dense} also uses
Section~\ref{sec:sse}.

\section{Preliminaries and the hardness input}\label{sec:prelim}

\subsection{Notation and conventions}

All graphs are finite. For a graph \(G\) and \(S\subseteq V(G)\), \(\partial
S=\partial_GS\) is the set of edges with exactly one endpoint in \(S\), and
\(e(S)\) is the number of edges with both endpoints in \(S\). If \(G\) is
\(d\)-regular and \(S\ne\varnothing\), the conductance of \(S\) is
\(\Phi_G(S)=|\partial S|/(d|S|)\).

A \emph{capacity system} \((V,c,\mu)\) consists of a finite set \(V\), a
symmetric function \(c:V\times V\to\Q_{\ge0}\) with \(c_{vv}=0\), and a
probability measure \(\mu\) on \(V\) with rational values. We write
\[
c(\partial S)=\sum_{u\in S,\,v\notin S}c_{uv},\qquad
W(c)=\sum_{\{u,v\}}c_{uv},\qquad d_v=\sum_uc_{uv},
\]
the first two sums being over pairs. Numbers are encoded in binary.

A \emph{fixed} quantity is chosen before the input formula and does not
depend on it. For a minimization problem, NP-hardness of approximation
within a factor \(C\) means that there is a deterministic polynomial-time
reduction mapping each 3CNF formula \(\varphi\) to an instance \(I\) and a
rational \(t>0\) with \(\mathrm{OPT}(I)\le t\) if \(\varphi\) is satisfiable
(YES) and \(\mathrm{OPT}(I)>Ct\) otherwise (NO); for maximization,
\(\mathrm{OPT}(I)\ge t\) in YES and \(\mathrm{OPT}(I)<t/C\) in NO. A
\(C\)-approximation algorithm would then decide 3SAT.

\subsection{Theorem U}

\begin{theoremU}[\cite{oai26}]
Fix a rational \(B\ge1\) and put \(A_0:=21\) and \(\tau:=1/20000\). There
are a constant \(W_0=W_0(B)\) and a deterministic polynomial-time reduction
mapping each 3CNF formula \(\varphi\) to a capacity system \((V,c,\mu)\) with
\(W(c)\le W_0\) and the following properties.
\begin{enumerate}[label=\textup{(U\arabic*)}]
\item If \(\varphi\) is satisfiable, there is \(S\subseteq V\) with
\(c(\partial S)\le A_0\) and \(\tau\le\mu(S)\le1-\tau\).
\item If \(\varphi\) is unsatisfiable, every \(S\subseteq V\) satisfies
\(c(\partial S)\ge B\mu(S)(1-\mu(S))\).
\end{enumerate}
\end{theoremU}

\paragraph{Reading Theorem U from \cite{oai26}.}
We record how Theorem~U follows from \cite{oai26}; all section and equation
numbers refer to that preprint. Its Theorem~6.1 provides, for each fixed
\(T_0\ge1\), a demand law \(\mu^{\rm raw}\) on scores and a finite sum \(D\)
of raw comparison forms of bounded total weight \(W^{\rm raw}(T_0)\), such
that at a perfect witness \(x\), available when \(\varphi\) is satisfiable,
the truth cost is at most \(4\) and
\(\mu^{\rm raw}\{B:B(x)\ge1/4\},\ \mu^{\rm raw}\{B:B(x)\le-1/4\}\ge10^{-4}\);
and such that, in a NO instance, every measurable binary lookup \(h\) with
\(\nu=\Var_{\mu^{\rm raw}}(h)>0\) satisfies \(D(h)>T_0\nu\).
Proposition~7.1 states that the seeds are polynomially enumerable, that the
numerical data range over a fixed finite family of polynomial patterns, and
that the total raw comparison mass is bounded by a constant. Section~8.1
rounds scores entrywise to a fixed grid, with error \(\delta_g\) satisfying
\(2\delta_gW^{\rm raw}\le1\) and \(\delta_g<1/16\), calls the rounded tables
keys, and computes in polynomial time (Lemma~8.1) rationals \(c_{ij}\) and
\(l_i\) with
\[
c_{ij}\ge c^0_{ij},\qquad 0\le l_i\le\pi_i,\qquad
\sum_{i<j}(c_{ij}-c^0_{ij})\le1,\qquad L:=\sum_il_i\ge1-10^{-6},
\]
where \(\pi\) is the pushforward of \(\mu^{\rm raw}\) to keys and
\(c^0_{ij}\) is the raw comparison weight between keys \(i\ne j\). Every
coloring of the keys extends to a measurable lookup, and
\(D(h)=\sum_{i<j}c^0_{ij}|h_i-h_j|\) (Equation~(122)).

Apply this with \(T_0=\max\{1,4B\}=4B\), stop at the end of Section~8.1
(the replication and coherence edges of Section~8.2 are not used), let
\(V\) be the set of keys, \(c\) the compiled capacities, and
\(\mu_i=l_i/L\). Then \(W(c)\le W^{\rm raw}+1=:W_0(B)\), and Equation~(125)
gives \(\mu_i\le2\pi_i\) and \(\mu(E)\ge\pi(E)-10^{-6}\) for every set
\(E\) of keys.

For (U2), let \(S\subseteq V\). If \(\mu(S)\mu(V\setminus S)=0\) there is
nothing to prove. Otherwise \(\pi(S)\pi(V\setminus S)\ge\mu(S)\mu(V\setminus
S)/4>0\), so the lookup \(h=\ind_S\) has positive raw variance and
\[
c(\partial S)\ \ge\ \sum_{i<j}c^0_{ij}|h_i-h_j|\ =\ D(h)\ >\ 4B\,\pi(S)\pi(V\setminus S)
\ \ge\ B\mu(S)(1-\mu(S)).
\]
For (U1), after rounding, the keys of the two truth tails have values
strictly above \(1/8\), respectively strictly below \(-1/8\), and each tail
retains \(\mu\)-mass at least \(10^{-4}-10^{-6}>\tau\). For each threshold
\(\theta\in[-1/8,1/8]\), color key \(i\) by \(\ind\{i(x)>\theta\}\); both
color classes have \(\mu\)-mass at least \(\tau\). For \(\theta\) uniform in
this interval of length \(1/4\), two values \(u,v\) are separated with
probability at most \(4|u-v|\). The rounded truth cost is at most \(5\), so
the expected cut capacity is at most \(4\cdot5\) plus the total excess
\(1\), that is, at most \(21=A_0\). Some threshold achieves this.

\begin{remark}
Section~8.1 of \cite{oai26} is written for the value of \(T_0\) fixed in
Section~8 there; its grid, key enumeration and Lemma~8.1 use only that
\(T_0\) is fixed, which is how we apply them.
\end{remark}

\subsection{Products and tensorization}

\begin{definition}
The \emph{product} of capacity systems \((V_j,c_j,\mu_j)\), \(j\in[s]\), is
the capacity system on \(V=\prod_jV_j\) with \(\mu=\bigotimes_j\mu_j\), in
which two points \(x,y\) that differ exactly in coordinate \(j\) have
capacity \(c_j(x_j,y_j)\prod_{i\ne j}\mu_i(x_i)\), and all other pairs have
capacity zero.
\end{definition}

Its total capacity is \(\sum_jW(c_j)\), since the masses of the other
coordinates sum to one.

\begin{lemma}[Tensorization]\label{lem:tensor}
If every factor satisfies \(c_j(\partial S)\ge B\mu_j(S)(1-\mu_j(S))\) for all
\(S\subseteq V_j\), then the product satisfies
\(c(\partial S)\ge B\mu(S)(1-\mu(S))\) for all \(S\subseteq V\).
\end{lemma}

\begin{proof}
Let \(f=\ind_S\). For \(x_{-j}\in\prod_{i\ne j}V_i\) put
\(S_{x_{-j}}=\{a\in V_j:(a,x_{-j})\in S\}\). Grouping the edges by their
moving coordinate,
\[
c(\partial S)=\sum_j\sum_{x_{-j}}\mu_{-j}(x_{-j})\,c_j(\partial S_{x_{-j}})
\ge B\sum_j\E\bigl[\Var(f\mid x_{-j})\bigr],
\]
since \(\mu_j(S_{x_{-j}})(1-\mu_j(S_{x_{-j}}))=\Var(f\mid x_{-j})\). Let
\(E_j\) be the averaging operator over coordinate \(j\) in \(L^2(\mu)\). The
\(E_j\) are commuting orthogonal projections, so \(f=\sum_{T\subseteq[s]}f_T\)
with \(f_T=\prod_{j\in T}(I-E_j)\prod_{j\notin T}E_jf\) is an orthogonal
decomposition, and \(\E\Var(f\mid x_{-j})=\|(I-E_j)f\|^2=\sum_{T\ni
j}\|f_T\|^2\). Summing over \(j\) gives
\(\sum_T|T|\|f_T\|^2\ge\sum_{T\ne\varnothing}\|f_T\|^2=\Var f=\mu(S)(1-\mu(S))\)
(this is the Efron--Stein inequality \cite{es81}).
\end{proof}

\subsection{Balanced colorings}

\begin{lemma}[Balanced colorings]\label{lem:balance}
Fix an integer \(K\ge2\) and a rational \(B\ge1\), and put
\(A_K:=A_0\tau^{-K}+1\). There are a constant \(W_1=W_1(K,B)\) and a
deterministic polynomial-time reduction mapping a 3CNF formula to a capacity
system \((V,c,\mu)\) together with an explicit action of the symmetric group
\(\mathfrak S_K\) on \(V\), such that the following hold.
\begin{enumerate}[label=\textup{(\roman*)}]
\item \(W(c)\le W_1\); \(c\) and \(\mu\) are invariant under the action; and
\(\mu(v)\le1/K\) for every \(v\in V\).
\item In YES there is a map \(\chi:V\to[K]\) with
\(\chi(\sigma v)=\sigma(\chi(v))\) for all \(\sigma\in\mathfrak S_K\), such
that the total capacity of pairs \(\{u,v\}\) with \(\chi(u)\ne\chi(v)\) is
at most \(A_K\). Consequently, for every \(\mathfrak S_K\)-invariant
\(\nu:V\to\Q_{\ge0}\), each class \(\chi^{-1}(k)\) has \(\nu\)-mass exactly
\(\nu(V)/K\).
\item In NO, \(c(\partial S)\ge B\mu(S)(1-\mu(S))\) for all \(S\subseteq V\).
\end{enumerate}
\end{lemma}

The constant \(A_K\) depends on \(K\) only, not on \(B\).

\begin{proof}
Let \((V_0,c_0,\mu_0)\) be the output of Theorem~U for this \(B\). Put
\(s_0=K\tau^K\) and fix \(t\) with \(B(1-s_0)^t\le1\); this choice uses only
the known bound \(s_0\), not a YES witness. Let \(F\) be the system on \([K]\)
with uniform measure and capacity \(B/K^2\) on every pair of distinct
symbols; for \(S\subseteq[K]\) it satisfies
\(c_F(\partial S)=|S|(K-|S|)B/K^2=B\mu_F(S)(1-\mu_F(S))\). Let \((V,c,\mu)\)
be the product of \(tK\) copies of \((V_0,c_0,\mu_0)\), indexed by
\((j,r)\in[t]\times[K]\) (block \(j\), position \(r\)), and one copy of
\(F\) (the fallback coordinate). The group \(\mathfrak S_K\) acts by
permuting the positions inside every block simultaneously and permuting the
fallback symbol by the same permutation; this preserves \(c\) and \(\mu\).
Item~(iii) follows from (U2) and Lemma~\ref{lem:tensor}. The total capacity
is \(tKW_0+B(K-1)/(2K)\le tKW_0+B=:W_1\), and the fallback coordinate gives
\(\mu(v)\le1/K\).

In YES let \(S_0\) be the set of (U1) and \(p=\mu_0(S_0)\in[\tau,1-\tau]\).
Call block \(j\) of \(x\in V\) \emph{successful} if exactly one position
\(r\) has \(x_{j,r}\in S_0\). Let \(\chi(x)\) be the \emph{position of the
unique coordinate in \(S_0\) inside the first successful block}, and let
\(\chi(x)\) be the fallback symbol if no block succeeds. Permuting positions
by \(\sigma\) preserves which blocks succeed and moves the distinguished
position by \(\sigma\), so \(\chi(\sigma x)=\sigma(\chi(x))\). As
\(\mathfrak S_K\) acts transitively on \([K]\) and \(\nu\) is invariant,
\(\nu(\chi^{-1}(\sigma k))=\nu(\sigma\chi^{-1}(k))=\nu(\chi^{-1}(k))\).

A block succeeds with probability \(s=Kp(1-p)^{K-1}\ge K\tau^K=s_0\). An
edge that changes coordinate \((j,r)\) can change \(\chi\) only if all
blocks before \(j\) fail and the two values of coordinate \((j,r)\) are
separated by \(S_0\). The event that the earlier blocks fail depends on the
other coordinates and has probability \((1-s)^{j-1}\), so these edges have
total capacity at most \(A_0(1-s)^{j-1}\) for each \((j,r)\). Edges changing
the fallback coordinate change \(\chi\) only if all \(t\) blocks fail. The
total capacity of pairs with different colors is therefore at most
\[
KA_0\sum_{j=1}^t(1-s)^{j-1}+\tfrac B2(1-s)^t\ \le\ \frac{KA_0}{s}+B(1-s_0)^t
\ \le\ A_0\tau^{-K}+1=A_K.\qedhere
\]
\end{proof}

\subsection{Finite-field incidence matrices}

Let \(q\) be a prime and \(\Omega=\F_q^2\). For \(t\in\F_q\) let
\(A_t\in\{0,1\}^{\Omega\times\Omega}\) be given by
\(A_t[(x,y),(u,v)]=1\) if and only if \(y+v-xu=t\), and for
\(T\subseteq\F_q\) put \(A_T=\sum_{t\in T}A_t\).

\begin{lemma}\label{lem:incidence}
\begin{enumerate}[label=\textup{(\roman*)}]
\item \(A_T\) is a symmetric \(0/1\) matrix with all row sums \(q|T|\);
different \(t\) give disjoint supports.
\item For every \(g\perp\ind\), \(\|A_tg\|\le\sqrt q\,\|g\|\) and
\(\|A_Tg\|\le|T|\sqrt q\,\|g\|\).
\item For \(S_1,S_2\subseteq\Omega\) of densities \(s_1,s_2\),
\[
\bigl|\ind_{S_1}^{\top}A_T\ind_{S_2}-q^3|T|s_1s_2\bigr|
\le|T|\,q^{5/2}\sqrt{s_1(1-s_1)s_2(1-s_2)}.
\]
\end{enumerate}
\end{lemma}

\begin{proof}
(i) The defining relation is symmetric in \(((x,y),(u,v))\); for fixed
\((x,y)\) and \(t\), the solutions are \(v=t-y+xu\), \(u\in\F_q\). Each
pair of points has a single value \(y+v-xu\).
(ii) The entry of \(A_tA_t^{\top}\) at \(((x,y),(x',y'))\) counts \((u,v)\)
with \(y+v-xu=t=y'+v-x'u\), that is \(y-y'=(x-x')u\). It is \(q\) on the
diagonal, \(0\) if \(x=x'\) and \(y\ne y'\), and \(1\) if \(x\ne x'\). Hence
\(A_tA_t^{\top}=qI+J-\bigoplus_xJ_q\), where \(J\) is the all-ones matrix
and \(\bigoplus_xJ_q\) the block matrix of all-ones blocks on the classes
of equal first coordinate. For \(g\perp\ind\) we get
\(\|A_t^{\top}g\|^2=q\|g\|^2-g^{\top}(\bigoplus_xJ_q)g\le q\|g\|^2\), and
\(A_t^{\top}=A_t\). The second bound follows by the triangle inequality.
(iii) Write \(\ind_{S_i}=s_i\ind+f_i\) with \(f_i\perp\ind\) and
\(\|f_i\|^2=q^2s_i(1-s_i)\). Since \(A_T\ind=q|T|\ind\), the cross terms
vanish and \(\ind^{\top}A_T\ind=q^3|T|\); the remaining term \(f_1^{\top}
A_Tf_2\) is bounded by (ii).
\end{proof}

We use \(A_T\) in two ways: as the biadjacency matrix between two disjoint
copies of \(\Omega\) (a simple bipartite graph), and, with its diagonal
removed, as the adjacency matrix of a simple graph on \(\Omega\). In the
second case the number of edges between \(S\) and \(\Omega\setminus S\) is
\(\ind_S^{\top}A_T\ind_{\Omega\setminus S}\), which the diagonal does not
affect. A prime in \([Q,2Q]\) exists by Bertrand's postulate and is found by
trial division in time polynomial in \(Q\).

\section{The treewidth family}\label{sec:tw}

\subsection{A common hard family}

Let \(\Tr=\Tr_{K,\ell}\) be the complete rooted \(K\)-ary tree of depth
\(\ell\), whose nodes are the words over \([K]\) of length at most \(\ell\)
and whose leaves are the words of length \(\ell\). For leaves \(a,b\), let
\(\route(a,b)\) be the set of nodes on the tree path from \(a\) to \(b\)
(so \(\route(a,a)=\{a\}\)).

\begin{theorem}[Common hard family]\label{th:family}
Fix an integer \(K\ge2\), a rational \(B\ge1\) and an integer \(n_1\). There
are a constant \(C_K\) depending only on \(K\), a constant
\(C'=C'(K,B)\), and a deterministic polynomial-time reduction mapping a
3CNF formula to a simple graph \(G\) on \(N\ge n_1\) vertices, a rational
\(\Lambda\ge N\) and an integer \(\ell\ge1\), with \(|E(G)|\le C'\Lambda\),
such that the following hold.
\begin{enumerate}[label=\textup{(Y\arabic*)}]
\item In YES, there is an ordering of \(V(G)\) every initial segment \(I\)
of which has \(|\partial I|\le C_K\Lambda\), and a partition of \(V(G)\)
into \(K\) classes of exactly \(N/K\) vertices with at most \(C_K\Lambda\)
edges joining different classes.
\item In YES, there is a map \(\lambda\) from \(V(G)\) to the leaves of
\(\Tr_{K,\ell}\) such that the loads
\(L_x=|\{uv\in E(G):x\in\route(\lambda(u),\lambda(v))\}|\) satisfy
\[
\max_{P}\sum_{x\in P}L_x\le C_K\Lambda,\qquad\sum_xL_x^2\le C_K\Lambda^2,
\]
the maximum being over root--leaf paths \(P\).
\end{enumerate}
\begin{enumerate}[label=\textup{(N)}]
\item In NO, every \(S\subseteq V(G)\) satisfies
\(|\partial S|\ge B\Lambda z(1-z)\), where \(z=|S|/N\).
\end{enumerate}
\end{theorem}

\subsection{Proof of Theorem~\ref{th:family}}

\paragraph{Step 1: balanced factor.}
Apply Lemma~\ref{lem:balance} with \(K\) and \(B\) and write
\((V_1,c_1,\mu_1)\), \(W_1\), \(\chi\) for its output.

\paragraph{Step 2: hierarchy.}
Fix \(\ell\ge1\) such that, with \(W:=\ell W_1\),
\begin{equation}\label{eq:ell}
WK^{-(\ell-1)}\le1,\qquad W^2K^{2-\ell}\le1,\qquad
\frac{B}{2K^\ell}\le\frac1{4(\ell+1)};
\end{equation}
this is possible since \(W\) grows linearly in \(\ell\). Let
\((V_2,c_2,\mu_2)\) be the product of \(\ell\) copies of
\((V_1,c_1,\mu_1)\). By Lemma~\ref{lem:tensor} it satisfies the NO
inequality, \(W(c_2)=W\), and every atom of \(\mu_2\) has mass at most
\(K^{-\ell}\). In YES put \(\lambda(v)=(\chi(v_1),\dots,\chi(v_\ell))\), a
leaf of \(\Tr\). By independence, for a word \(p\) of length \(j\) the set of
\(v\) with \(\lambda(v)\) beginning with \(p\) has \(\mu_2\)-mass exactly
\(K^{-j}\).

An edge of \(c_2\) changes one coordinate \(i\); it changes at most the
label \(\lambda_i\). We use the following count: \emph{the capacity of edges
that change coordinate \(i\) and the label \(\lambda_i\), and whose
endpoints have prescribed labels in \(j'\) other coordinates, is at most
\(A_KK^{-j'}\)}, because the label-changing capacity in one factor is at most
\(A_K\) and the prescribed labels of the other coordinates have probability
\(K^{-j'}\). Likewise edges that change coordinate \(i\) but not
\(\lambda_i\), with all other \(\ell-1\) labels prescribed, have capacity
at most \(W_1K^{-(\ell-1)}\); summing over \(i\), the capacity of edges with
both endpoints in one leaf is at most \(WK^{-(\ell-1)}\le1\).

\begin{claim}\label{cl:weighted}
In YES:
\begin{enumerate}[label=\textup{(\alph*)}]
\item the partition of \(V_2\) by \(\lambda_1\) has classes of
\(\mu_2\)-mass \(1/K\) and crossing capacity at most \(A_K\);
\item for every ordering of \(V_2\) that is lexicographic in \(\lambda\)
(ties arbitrary), every initial segment has capacity at most \(2A_K+1\);
\item the loads \(L^{(2)}_x=\sum c_2(u,v)\), over pairs with
\(x\in\route(\lambda(u),\lambda(v))\), satisfy for a node \(x\) of depth
\(j\)
\[
L^{(2)}_x\le(Kj+1)A_KK^{-j}+\ind_{\{j=\ell\}}WK^{-(\ell-1)} .
\]
Consequently \(\sum_{x\in P}L^{(2)}_x\le P_K:=A_K\sum_{j\ge0}(Kj+1)K^{-j}+1\)
for every root--leaf path \(P\), and \(\sum_x(L^{(2)}_x)^2\le
Q_K:=2A_K^2\sum_{j\ge0}(Kj+1)^2K^{-j}+2\).
\end{enumerate}
\end{claim}

\begin{proof}
(a) is the case \(j'=0\), \(i=1\) of the count. (b) Let \(I\) be an initial
segment and \(w\) the label of the first vertex outside \(I\). An edge
\(\{u,v\}\) crossing \(I\) with \(\lambda(u)\ne\lambda(v)\) changes exactly
one label \(i\), so \(\lambda(u),\lambda(v)\) share a prefix \(p\) of length
\(i-1\). If \(p\) differs from the first \(i-1\) letters of \(w\), the whole
subtree of \(p\) lies on one side of \(I\); so \(p=w_{<i}\), and these edges
have capacity at most \(A_KK^{-(i-1)}\). Edges with
\(\lambda(u)=\lambda(v)\) can cross only inside the leaf \(w\), with capacity
at most \(1\). Summing, the cut is at most \(A_K\sum_{i\ge1}K^{-(i-1)}+1\le
2A_K+1\).
(c) Let \(x\) have depth \(j\) and word \(p\). A route through \(x\) either
has its splitting node at \(x\), or has exactly one endpoint in the subtree
of \(x\), or (if \(j=\ell\)) lies inside the leaf \(x\). Edges splitting at
\(x\) change label \(j+1\) and have prescribed labels \(p\) in \(j\) other
coordinates: capacity at most \(A_KK^{-j}\). An edge entering the subtree
from outside changes one label \(i\le j\), and its inner endpoint has the
labels of \(p\) in the other \(j-1\) of the first \(j\) coordinates:
capacity at most \(A_KK^{-(j-1)}\) for each \(i\), in total
\(KjA_KK^{-j}\). Leaf-internal edges contribute at most \(WK^{-(\ell-1)}\).
Summing along a path gives \(P_K\) since \(WK^{-(\ell-1)}\le1\). There are
\(K^j\) nodes at depth \(j\), so
\(\sum_x(L^{(2)}_x)^2\le\sum_jK^j\cdot2(Kj+1)^2A_K^2K^{-2j}+2K^\ell W^2
K^{-2(\ell-1)}\le Q_K\), using \(W^2K^{2-\ell}\le1\).
\end{proof}

\paragraph{Step 3: exact vertex counts.}
Put \(\eta=1/(\ell+1)\) and let \(h=|V_2|\). For an integer \(J\ge1\) let
\(n_v=K\max\{1,\lceil J\mu_2(v)\rceil\}\) for \(v\in V_2\), replace \(v\) by
a cloud of \(n_v\) copies, and let \(V_3\) be the union of the clouds,
\(n'=|V_3|\), \(\rho\) the uniform measure on \(V_3\), \(\nu(a)=\mu_2(v)/n_v\)
for \(a\) in the cloud of \(v\), and \(\omega=KJ/n'\). Then
\(KJ\le n'\le K(J+h)\), \(K\mid n'\), \(1-\omega\le h/J\), and
\(\rho\ge\omega\nu\) pointwise. Define capacities on \(V_3\):
\begin{align*}
c^{\rm inh}(a,b)&=c_2(v,w)/(n_vn_w)&&\text{for \(a,b\) in the clouds of \(v\ne w\)};\\
c^{\rm cl}(a,b)&=B\nu(a)\nu(b)&&\text{for distinct \(a,b\) in the same cloud};\\
c^{\rm cor}(a,b)&=B\bigl(\rho(a)\rho(b)-\omega^2\nu(a)\nu(b)\bigr)\ \ge0
&&\text{for all distinct \(a,b\)};
\end{align*}
and \(c_3=c^{\rm inh}+c^{\rm ext}\) with \(c^{\rm ext}=c^{\rm cl}+c^{\rm cor}\).

\begin{claim}\label{cl:counts}
\begin{enumerate}[label=\textup{(\alph*)}]
\item In NO, \(c_3(\partial S)\ge B\rho(S)(1-\rho(S))\) for every
\(S\subseteq V_3\).
\item \(W(c^{\rm ext})\le B/(2K^\ell)+B(1-\omega^2)/2\).
\item Suppose \(J\) is so large that \(B(1-\omega^2)/2\le\eta/4\) and
\(2K(1-\omega)(W+1)\le1\). In YES there is a partition of \(V_3\) into \(K\)
classes of exactly \(n'/K\) elements whose crossing \(c_3\)-capacity is at
most \(A_K+2\).
\end{enumerate}
\end{claim}

\begin{proof}
(a) Let \(a_v\) be the fraction of the cloud of \(v\) lying in \(S\). Let
\(R\subseteq V_2\) contain each \(v\) independently with probability
\(a_v\). Since capacities are spread uniformly over pairs of clouds,
\[
c^{\rm inh}(\partial S)=\sum_{\{v,w\}}c_2(v,w)\bigl[a_v(1-a_w)+a_w(1-a_v)\bigr]
=\E\,c_2(\partial R)\ge B\,\E\bigl[\mu_2(R)(1-\mu_2(R))\bigr].
\]
With \(z_\nu=\sum_v\mu_2(v)a_v=\nu(S)\), the last expectation equals
\(z_\nu(1-z_\nu)-\sum_v\mu_2(v)^2a_v(1-a_v)\). The cloud edges contribute
\(c^{\rm cl}(\partial S)=\sum_vB(\mu_2(v)/n_v)^2(a_vn_v)((1-a_v)n_v)
=B\sum_v\mu_2(v)^2a_v(1-a_v)\). Hence
\((c^{\rm inh}+c^{\rm cl})(\partial S)\ge B\nu(S)\nu(V_3\setminus S)
=B\sum_{a\in S,b\notin S}\nu(a)\nu(b)\ge B\omega^2\sum_{a\in S,b\notin
S}\nu(a)\nu(b)\), and adding \(c^{\rm cor}(\partial S)\) gives at least
\(B\sum_{a\in S,b\notin S}\rho(a)\rho(b)=B\rho(S)(1-\rho(S))\).

(b) \(W(c^{\rm cl})\le\frac B2\sum_v\mu_2(v)^2\le\frac B2\max_v\mu_2(v)\le
B/(2K^\ell)\). Next,
\(W(c^{\rm cor})=\frac B2\bigl[(1-\sum\rho^2)-\omega^2(1-\sum\nu^2)\bigr]\),
and \(\omega^2\sum_a\nu(a)^2\le\sum_a\rho(a)^2\) since \(\omega\nu\le\rho\);
so \(W(c^{\rm cor})\le B(1-\omega^2)/2\).

(c) Consider the classes \(\Gamma_k=\bigcup\{\text{cloud of }v:\chi(v_1)=k\}\).
Since \(n_v\ge KJ\mu_2(v)\) and the \(\lambda_1\)-classes have
\(\mu_2\)-mass \(1/K\), each \(|\Gamma_k|\ge J=\omega n'/K\). Thus the total
excess \(\sum_k(|\Gamma_k|-n'/K)_+\), which equals the total deficit, is at
most \((1-\omega)n'\). By (b), \eqref{eq:ell} and the choice of \(J\),
\(W(c^{\rm ext})\le\eta/2\le1\), so the \(c_3\)-degrees sum to at most
\(2(W+1)\). From each oversize class move its surplus vertices of smallest
\(c_3\)-degree to deficient classes. An oversize class has at least \(n'/K\)
members, so its \(m\) smallest degrees sum to at most \(2m(W+1)K/n'\); the
moves increase the crossing capacity by at most \(2K(1-\omega)(W+1)\le1\).
Before the moves, the crossing capacity is at most \(A_K\) (inherited part,
Claim~\ref{cl:weighted}(a)) plus \(\eta/2\).
\end{proof}

In YES give each element of the cloud of \(v\) the leaf \(\lambda(v)\). For
every node \(x\), the \(c^{\rm inh}\)-load equals \(L^{(2)}_x\). Order
\(V_3\) lexicographically by label, with each cloud contiguous. An initial
segment splits at most one cloud, and its \(c^{\rm inh}\)-capacity is
linear in the fraction of that cloud; hence it is at most the larger of the
two adjacent whole-cloud segments, which are initial segments of a
lexicographic order of \(V_2\), so at most \(2A_K+1\) by
Claim~\ref{cl:weighted}(b).

Choose \(J\), polynomial in the input, so large that the hypotheses of
Claim~\ref{cl:counts}(c) hold and \(B/n'\le\eta/8\); this is possible since
\(1-\omega\le h/J\) and \(n'\ge KJ\). Then
\(W(c^{\rm ext})\le\eta/4+\eta/4=\eta/2\) by \eqref{eq:ell}.

\paragraph{Step 4: a simple graph.}
Let \(U\) be an integer with \(U>W+1\ge W(c_3)\ge\max c_3\). Choose a prime
\(q\ge16\), polynomial in \(n'\), such that
\[
q^2\ge Un',\qquad n'q^2\ge n_1,\qquad
\frac2U\Bigl(\frac{B}{(n')^2}+\frac{Un'}{\sqrt q}\Bigr)+\frac1q\le1,\qquad
\frac{U(n')^2}{\sqrt q}+\frac{Un'}{2q}+\frac{U(n')^2}{2q}\le\frac\eta4 .
\]
Put \(M=q^2\) and \(\Lambda=M^2/U\). The graph \(G\) has vertex set
\(V_3\times\F_q^2\), with clouds \(C_a=\{a\}\times\F_q^2\) and \(N=n'q^2\).
For distinct \(a,b\in V_3\) choose \(T_{ab}\subseteq\F_q\) with
\(|T_{ab}|=\lceil qc_3(a,b)/U\rceil\le q\) and join \(C_a\) to \(C_b\) by
\(A_{T_{ab}}\). Inside each \(C_a\) take \(A_{T_a}\) without its diagonal,
where \(|T_a|=\bigl\lceil q\cdot\frac2U\bigl(\frac{B}{(n')^2}+\frac{Un'}{\sqrt
q}\bigr)\bigr\rceil\le q\). By Lemma~\ref{lem:incidence}(i), \(G\) is simple,
and \(\Lambda\ge N\) because \(q^2\ge Un'\).

\emph{NO.} Let \(S\subseteq V(G)\), let \(a_w\) be the density of \(S\) in
\(C_w\), \(u_w=\sqrt{a_w(1-a_w)}\) and \(z=|S|/N=\frac1{n'}\sum_wa_w\). For a
pair \(\{w,w'\}\) the crossing edges are
\(\ind_{S_w}^{\top}A\ind_{C_{w'}\setminus S_{w'}}+\ind_{C_w\setminus
S_w}^{\top}A\ind_{S_{w'}}\) with \(A=A_{T_{ww'}}\). By
Lemma~\ref{lem:incidence}(iii), \(qc_3(w,w')/U\le|T_{ww'}|\le q\) and
\(\Lambda=q^4/U\), this is at least
\(\Lambda\bigl(c_3(w,w')[a_w(1-a_{w'})+a_{w'}(1-a_w)]-\frac{2U}{\sqrt q}u_wu_{w'}\bigr)\).
Summing over pairs and using \(\sum_{w<w'}u_wu_{w'}\le\frac{n'}2\sum_wu_w^2\),
the intercloud edges number at least
\[
\Lambda\Bigl(\E\,c_3(\partial R)-\frac{Un'}{\sqrt q}\sum_wu_w^2\Bigr)
\ \ge\ \Lambda\Bigl(Bz(1-z)-\Bigl(\frac{B}{(n')^2}+\frac{Un'}{\sqrt q}\Bigr)\sum_wu_w^2\Bigr),
\]
where \(R\subseteq V_3\) contains \(w\) independently with probability
\(a_w\), and we used Claim~\ref{cl:counts}(a) and
\(\E[\rho(R)(1-\rho(R))]=z(1-z)-\sum_wu_w^2/(n')^2\). Inside \(C_w\),
Lemma~\ref{lem:incidence}(iii) gives at least
\(q^3|T_w|u_w^2(1-q^{-1/2})\ge\Lambda\,U\frac{|T_w|}{q}\cdot\frac12u_w^2
\ge\Lambda\bigl(\frac{B}{(n')^2}+\frac{Un'}{\sqrt q}\bigr)u_w^2\) crossing
edges. Adding, \(|\partial S|\ge B\Lambda z(1-z)\).

\emph{Edge counts.} Between \(C_a\) and \(C_b\) there are
\(q^3|T_{ab}|=\Lambda U|T_{ab}|/q\in[\Lambda c_3(a,b),\Lambda(c_3(a,b)+U/q)]\)
edges, and inside \(C_a\) at most \(q^3|T_a|/2\le\Lambda\bigl(\frac B{(n')^2}
+\frac{Un'}{\sqrt q}+\frac U{2q}\bigr)\). Call \emph{extra} all edges beyond
the \(\Lambda c^{\rm inh}\) prescribed by the inherited capacities: their
number divided by \(\Lambda\) is at most
\(W(c^{\rm ext})+\frac{U(n')^2}{2q}+\frac B{n'}+\frac{U(n')^2}{\sqrt q}
+\frac{Un'}{2q}\le\frac\eta2+\frac\eta8+\frac\eta4\le\eta\). In particular
\(|E(G)|\le(W+2)\Lambda=:C'\Lambda\).

\emph{YES partition.} Lifting the partition of Claim~\ref{cl:counts}(c) by
whole clouds gives \(K\) classes of exactly \(N/K\) vertices with at most
\(\Lambda(A_K+2+\eta)\le\Lambda(A_K+3)\) crossing edges.

\emph{YES order.} List the clouds in the order of \(V_3\) chosen in Step~3,
each cloud in any internal order. An initial segment \(I\) consists of the
whole clouds of an initial segment \(I_3\) of \(V_3\) and part of one cloud
\(C_{a^*}\). Only one cloud is split, so the vectors \(f_w\) of
Lemma~\ref{lem:incidence}(iii) vanish for all other clouds and the
intercloud term has no discrepancy: the intercloud crossing edges number
\(\Lambda\sum_{w<w'}(U|T_{ww'}|/q)[a_w(1-a_{w'})+a_{w'}(1-a_w)]\), which is
linear in \(a_{a^*}\) and hence at most its value at \(a_{a^*}\in\{0,1\}\).
That value is at most \(\Lambda\) times the \(c^{\rm inh}\)-capacity of an
initial segment of \(V_3\) plus the extra edges, so \(|\partial I|\le
\Lambda(2A_K+1+\eta)\le\Lambda(2A_K+2)\).

\emph{YES routes.} Give each vertex of \(C_a\) the leaf of \(a\). Then
\(L_x\le\Lambda(L^{(2)}_x+e_x)\), where \(e_x\) is the load of the extra
edges, divided by \(\Lambda\). Extra edges have total scaled count at most
\(\eta\), and each route meets at most two nodes of each depth, so the
nodes of one depth carry total extra load at most \(2\eta\). Along a path,
\(\sum_{x\in P}e_x\le(\ell+1)\eta=1\), and
\(\sum_xe_x^2\le(\ell+1)(2\eta)^2\le4\). Therefore
\(\sum_{x\in P}L_x\le\Lambda(P_K+1)\) and
\(\sum_xL_x^2\le\Lambda^2(2Q_K+8)\).

Taking \(C_K=\max\{2A_K+3,\,P_K+1,\,2Q_K+8\}\), which depends only on \(K\),
proves the theorem. All steps are polynomial for fixed \(K,B,n_1\). \qed

\subsection{Edge cuts, layouts and partitioning}\label{sec:edge}

Throughout, \(G,\Lambda,N\) come from Theorem~\ref{th:family}; \(B\) is
taken large after \(K\) and the target factor are fixed.

\begin{corollary}\label{cor:edge}
Assume Theorem~U. For every fixed \(C\ge1\) the following problems are
NP-hard to approximate within \(C\):
\begin{enumerate}[label=\textup{(\alph*)}]
\item minimum bisection; for fixed \(0<\alpha\le1/2\), the minimum of
\(|\partial S|\) over \(\alpha N\le|S|\le(1-\alpha)N\); for fixed
\(1/2<\theta<1\), the minimum number of edges whose deletion leaves only
components of at most \(\theta N\) vertices;
\item cutwidth; minimum linear arrangement; every fixed \(\ell_p\)-norm
\((1\le p\le\infty)\) of the vector of prefix cut sizes of an ordering;
carving-width;
\item for fixed \(k\ge2\) and \(1\le\beta<k\), min-sum and min-max
balanced \(k\)-partitioning, where YES has a partition into parts of exactly
\(N/k\) vertices and NO bounds every partition into \(k\) parts of at most
\(\beta N/k\) vertices;
\item for fixed rational \(0<\delta<1\) and fixed \(\gamma\) with
\(\gamma\delta<1\), the minimum of \(|\partial S|/|S|\): YES has a set of
exactly \(\delta N\) vertices with \(|\partial S|/|S|\le C_K\Lambda/(\delta
N)\), and NO has \(|\partial S|/|S|\ge B(1-\gamma\delta)\Lambda/N\) for
every nonempty \(S\) with \(|S|\le\gamma\delta N\).
\end{enumerate}
\end{corollary}

\begin{proof}
(a) Take \(K\) even. The union of \(K/2\) classes of (Y1) is a bisection
with at most \(C_K\Lambda\) crossing edges. In NO, every \(S\) with
\(\alpha\le|S|/N\le1-\alpha\) has \(|\partial S|\ge B\alpha(1-\alpha)\Lambda\).
For the component version, the YES bisection's edges leave components of at
most \(N/2\le\theta N\) vertices. In NO, let \(F\) be such an edge set and
\(a=(1-\theta)/2\). If a component has at least \(aN\) vertices, take it as
\(S\); its complement has at least \((1-\theta)N\ge aN\) vertices. Otherwise
accumulate components until their total first reaches \(aN\); the total is
below \(2aN=(1-\theta)N\), so the complement exceeds \(\theta N\ge aN\).
Every edge of \(\partial S\) joins different components, so lies in \(F\),
and \(|F|\ge B\Lambda a(1-a)\).

(b) The YES order of (Y1) has every prefix cut at most \(C_K\Lambda\), so
cutwidth at most \(C_K\Lambda\), minimum linear arrangement (the sum of the
prefix cuts) at most \(NC_K\Lambda\), and \(\ell_p\)-norm at most
\(N^{1/p}C_K\Lambda\). In NO, every ordering has at least \(N/2-1\) prefixes
of size in \([N/4,3N/4]\), each with cut at least \(3B\Lambda/16\); this
gives the lower bounds \(3B\Lambda/16\), \((N/2-1)3B\Lambda/16\) and
\((N/2-1)^{1/p}3B\Lambda/16\). For carving-width (decompositions by
subcubic trees whose leaves are the vertices), every such tree has an edge
whose two sides contain at least \((N-1)/3\) vertices each: at a node \(c\)
all of whose branches contain at most \(N/2\) leaves, the at most three
branches contain at least \(N-1\) leaves, and the largest branch gives the
edge. So the NO carving-width is at least \(B\Lambda z(1-z)\) with
\(z\in[(N-1)/(3N),1/2]\). In YES use a caterpillar along the YES order: its
nonsingleton cuts are prefix cuts, and a singleton cut \(\{v\}\) costs
\(\deg v\le|\partial I^-|+|\partial I^+|\le2C_K\Lambda\), where \(I^-\) and
\(I^+\) are the prefixes ending just before and at \(v\).

(c) Take \(K\) divisible by \(k\) and group the YES classes into \(k\) parts
of \(N/k\) vertices; each part has boundary at most \(C_K\Lambda\), so the
sum of boundaries is at most \(kC_K\Lambda\). In NO let \(z_1,\dots,z_k\le
\beta/k\) be the part masses; then
\(\sum_i|\partial S_i|\ge B\Lambda\sum_iz_i(1-z_i)\ge B\Lambda(1-\beta/k)\),
and the maximum is at least \(1/k\) of the sum.

(d) Take \(K\) divisible by the denominator of \(\delta\); the union of
\(K\delta\) YES classes has exactly \(\delta N\) vertices and boundary at
most \(C_K\Lambda\). In NO, \(|\partial S|\ge B\Lambda z(1-z)\ge
B\Lambda z(1-\gamma\delta)\) for \(z\le\gamma\delta\).

In each case the ratio between the NO lower bound and the YES upper bound
is a fixed positive multiple of \(B/C_K\) (asymptotically in \(N\), which
Theorem~\ref{th:family} lets us take large), so a large fixed \(B\) gives
factor \(C\).
\end{proof}

\subsection{Vertex separators, width parameters and fill-in}\label{sec:vertex}

\paragraph{The incidence graph.}
Following \cite{apw12,wapl14}, for a graph \(G\) and an integer \(r\ge1\)
let \(H_r(G)\) have, for every \(v\in V(G)\), \(r\) \emph{copies} of \(v\),
and for every edge \(uv\) one \emph{edge vertex} adjacent to all copies of
\(u\) and of \(v\). It is bipartite with \(rN+|E(G)|\) vertices. Let \(G\)
come from Theorem~\ref{th:family} with \(K\) even, put \(m=|E(G)|\),
\(r=100N(m+1)\), \(H=H_r(G)\) and
\[
T=B\Lambda/16 .
\]
In NO, every cut of \(G\) with both sides of at least \(N/8\) vertices has
at least \(B\Lambda\cdot\frac18\cdot\frac78\ge T\) edges; in particular
\(T\le m\).

A set \(X\subseteq V(H)\) is a \emph{\(\theta\)-balanced separator} if every
component of \(H-X\) has at most \(\theta|V(H)|\) vertices. We use the
standard fact that every tree decomposition of width \(w\) has a bag that
is a \(\tfrac12\)-balanced separator: assign each vertex to a bag containing
it, choose a node \(t^*\) such that every component of the tree minus
\(t^*\) receives at most \(|V(H)|/2\) assigned vertices, and note that the
bags meeting a component of \(H-B_{t^*}\) form a subtree avoiding \(t^*\).
Hence \(\tw(H)+1\) is at least the minimum size of a \(\tfrac23\)-balanced
separator.

\begin{lemma}[NO bounds]\label{lem:sepNO}
In NO, every \(\tfrac23\)-balanced separator of \(H\) has at least \(T\)
vertices. Consequently \(\tw(H)\ge T-1\), \(\pw(H)\ge T-1\) and
\(\td(H)\ge T\).
\end{lemma}

\begin{proof}
Let \(|X|<T\). Let \(F\) be the set of edges of \(G\) whose edge vertices lie
in \(X\), and \(G'=G-F\). Since \(|X|<T\le m\), every \(v\in V(G)\) keeps at
least \(r-m\) copies. For a component \(Q\) of \(G'\) with \(|Q|\ge2\), the
surviving copies of the vertices of \(Q\) and the surviving edge vertices of
edges of \(G'[Q]\) induce a connected subgraph of \(H-X\). If some component
has \(|Q|>3N/4\), then \(H-X\) has a component with more than
\(\frac{3N}4(r-m)>\frac23(rN+m)\) vertices, by the choice of \(r\). Otherwise
group the components of \(G'\) into \(S\) with \(N/8\le|S|\le7N/8\) as in
the proof of Corollary~\ref{cor:edge}(a); all edges of \(\partial_GS\) lie
in \(F\), so \(|X|\ge|F|\ge T\), a contradiction. The bound on \(\tw\)
follows from the fact above; \(\pw\ge\tw\); and an elimination forest of
height \(h\) yields a tree decomposition whose bags are the root-to-vertex
chains, so \(\td\ge\tw+1\).
\end{proof}

\begin{lemma}[YES bounds]\label{lem:sepYES}
In YES:
\begin{enumerate}[label=\textup{(\alph*)}]
\item \(H\) has a \(\tfrac23\)-balanced separator of at most \(C_K\Lambda\)
vertices;
\item there is an ordering of \(V(H)\) in which, at every prefix, at most
\(2C_K\Lambda\) placed vertices have an unplaced neighbor; hence
\(\pw(H)\le2C_K\Lambda\);
\item \(\td(H)\le C_K\Lambda+1\);
\item \(H\) has a chordal supergraph with at most \(C_K\Lambda^2/2\) added
edges, none of them incident to a copy.
\end{enumerate}
\end{lemma}

\begin{proof}
(a) Let \(S\) be the union of \(K/2\) classes of (Y1) and \(X\) the edge
vertices of \(\partial_GS\). Each component of \(H-X\) has at most
\(rN/2+m\le\frac23(rN+m)\) vertices.

(b) Let \(v_1,\dots,v_N\) be the YES order of (Y1). For \(i=1,\dots,N\)
place the edge vertices of the edges from \(v_i\) to later vertices, then
the \(r\) copies of \(v_i\). A copy is never live (all its neighbors are
edge vertices placed earlier). Between groups, the live vertices are the
edge vertices of \(\partial\{v_1,\dots,v_i\}\); during group \(i\) they lie
in \(\partial\{v_1,\dots,v_{i-1}\}\cup\partial\{v_1,\dots,v_i\}\). So at
most \(2C_K\Lambda\) vertices are live at any time. Taking as bags the
current vertex together with the vertices live just before it gives a path
decomposition of width at most \(2C_K\Lambda\): an edge is covered in the bag
of its later endpoint, and each vertex occupies the consecutive bags from its
placement to its last neighbor.

(c) Use the map \(\lambda\) of (Y2). For a node \(x\) let \(R_x\) be the set
of edge vertices \(uv\) with \(x\) the splitting node of
\(\lambda(u),\lambda(v)\) (for a leaf: \(\lambda(u)=\lambda(v)=x\)). Build a
rooted forest recursively: at \(x\), stack \(R_x\) as a chain and hang
below its last element the forests built at the children of \(x\); at a
leaf, hang the copies of the vertices assigned to it below the chain. Every
edge of \(H\) joins an edge vertex \(e=uv\in R_x\) to a copy of \(u\) or
\(v\), which lies below \(x\) and hence below \(e\). Since \(|R_x|\le L_x\),
the height is at most \(\max_P\sum_{x\in P}L_x+1\le C_K\Lambda+1\).

(d) Add to \(\Tr\) a private child \(y_\kappa\) of the leaf \(\lambda(v)\)
for each copy \(\kappa\) of each \(v\). Represent the copy \(\kappa\) by
\(\{y_\kappa\}\) and the edge vertex \(e=uv\) by the subtree
\(\route(\lambda(u),\lambda(v))\cup\{y_\kappa:\kappa\text{ a copy of }u\text{ or
}v\}\). The intersection graph \(H'\) of these subtrees contains \(H\). It is
chordal: each subtree has a highest node; eliminate repeatedly a subtree
whose highest node \(h_a\) is deepest. A remaining neighbor meets it below
or at \(h_a\) and has its own highest node at depth at most that of \(h_a\),
so it contains \(h_a\); the remaining neighbors pairwise intersect at
\(h_a\). In \(H'\) copies are adjacent only to the edge vertices of their
incident edges, and two edge-vertex subtrees meeting at a private child also
meet at its parent. So the added edges join edge vertices whose routes share
a node of \(\Tr\), and their number is at most
\(\sum_x\binom{L_x}2\le\frac12\sum_xL_x^2\le C_K\Lambda^2/2\).
\end{proof}

\begin{proposition}[Fill-in]\label{prop:fillin}
In NO, every chordal supergraph of \(H\) has at least \(T^2/8\) added edges.
\end{proposition}

\begin{proof}
A chordal supergraph \(H'\) has treewidth at least \(\tw(H)\ge T-1\), so it
has a clique on \(s\ge T\) vertices. As \(H\) is bipartite, at most
\(s^2/4\) of its edges lie inside this clique, so at least
\(\binom s2-s^2/4=s(s-2)/4\ge s^2/8\) edges are added (\(s\ge4\)).
\end{proof}

\paragraph{Other width parameters.}
We use the following comparisons, for every graph \(H\):
\(\mmw(H)\le\bw(H)\le\tw(H)+1\le3\mmw(H)\) \cite{rs91,vat12,jst15}; on the
line graph \(L(H)\), clique-width lies between \((\tw(H)+1)/4\) and
\(2\tw(H)+2\) \cite{gw07}, and NLC-width is within a factor two of
clique-width \cite{joh98}. Branch decompositions of \(H\) are subcubic trees
with leaf set \(E(H)=V(L(H))\), which are also the decomposition trees for
rank-width, Boolean-width and mim-width of \(L(H)\). For a tree edge
splitting \(E(H)\) into \(E_1,E_2\), let \(h\) be the number of vertices of
\(H\) incident to edges of both sides (the middle set; \(\bw\) minimizes the
maximum of \(h\)).

\begin{lemma}\label{lem:linewidth}
For a bipartite simple graph \(H\) and every cut \((E_1,E_2)\) as above, the
mim-value, the \(\F_2\)-cut-rank and the Boolean-width value
(\(\log_2\) of the number of distinct neighborhood unions) of the cut in
\(L(H)\) all lie between \(h/8\) and \(h\). Hence rank-width, Boolean-width
and mim-width of \(L(H)\) lie between \(\bw(H)/8\) and \(\bw(H)\).
\end{lemma}

\begin{proof}
Upper bounds: \(e_1\in E_1\) and \(e_2\in E_2\) are adjacent in \(L(H)\) iff
they share a vertex, which is then a middle vertex. The cross matrix equals
\(XY^{\top}\), where \(X,Y\) are the incidence matrices of \(E_1,E_2\)
against the middle set; over \(\F_2\) this is exact because two distinct
edges share at most one endpoint. So the cut-rank is at most \(h\). A union
of neighborhoods of edges of \(E_1\) is determined by the middle vertices
they touch, so there are at most \(2^h\) such unions. An induced matching
across the cut gives an identity submatrix of the cross matrix and \(2^s\)
distinct unions, so mim-value \(\le\) cut-rank and mim-value \(\le\)
Boolean-width value; it remains to show that the mim-value is at least
\(h/8\). At least \(h/2\) middle vertices lie in one class of the
bipartition of \(H\); at each such \(v\) choose edges \(vl_v\in E_1\) and
\(vr_v\in E_2\), where \(l_v\ne r_v\) lie in the other class. Color that
other class randomly red/blue; in expectation at least \(h/8\) of these
\(v\) have \(l_v\) red and \(r_v\) blue. For such \(v\ne v'\), the edges
\(vl_v\) and \(v'r_{v'}\) share no vertex (\(v\ne v'\) lie in the same
class, and \(l_v\ne r_{v'}\) by color), so the pairs
\((vl_v,vr_v)\) form an induced matching across the cut.
\end{proof}

The lower bound \(\bw(H)/25\) for mim-width of line graphs from
\cite[Theorem~4]{bmpy} would also suffice.

\begin{proposition}[Tree-independence number]\label{prop:treealpha}
Let \(\mathrm{tree}\text{-}\alpha\) denote the tree-independence number. In
YES, \(\mathrm{tree}\text{-}\alpha(L(H))\le\tw(H)+1\le2C_K\Lambda+1\). In NO,
\(\mathrm{tree}\text{-}\alpha(L(H))\ge T/2\).
\end{proposition}

\begin{proof}
YES: replace each bag of a tree decomposition of \(H\) by the set of edges
incident to its vertices. The bags containing \(uv\) form the union of the
subtrees of \(u\) and \(v\), which intersect; adjacent edges share a vertex
and so share a bag. An independent set in a new bag is a matching meeting
the old bag in distinct vertices, so it has at most \(\tw(H)+1\) elements;
and \(\tw(H)\le\pw(H)\le2C_K\Lambda\). NO: take a tree decomposition of
\(L(H)\) whose bags have independence number at most \(a\). The edges at a
vertex \(v\) of \(H\) form a clique of \(L(H)\), so some bag contains them
all; assign \(v\) to such a bag (isolated vertices arbitrarily). Choose a
node \(t^*\) such that each component of the tree minus \(t^*\) receives at
most \(|V(H)|/2\) assigned vertices, and let \(Y\) be the endpoints of a
maximal matching among the edges of the bag at \(t^*\); then \(|Y|\le2a\) and
\(Y\) covers every edge in that bag. An edge \(uv\) whose endpoints are
assigned to different components, or one of them to \(t^*\), lies in the
bags of \(u\) and \(v\) and hence, by connectivity, in the bag at \(t^*\);
so it is covered by \(Y\). Surviving vertices assigned to \(t^*\) are
isolated in \(H-Y\). Thus \(Y\) is a \(\frac12\)-balanced separator, and
\(2a\ge T\) by Lemma~\ref{lem:sepNO}.
\end{proof}

\subsection{Vertex layouts, interval completion and pebbling}\label{sec:layouts}

For an ordering \(\pi\) of a graph and \(0\le i\le n\), let \(\live_\pi(i)\)
be the number of vertices among the first \(i\) that have a neighbor after
position \(i\). The \emph{vertex separation} is \(\max_i\live_\pi(i)\)
(equal to pathwidth \cite{kin92}) and the \emph{vertex sum} is
\(\sum_i\live_\pi(i)\); cutwidth and minimum linear arrangement are the
edge analogues. For a DAG and a topological order, the directed live count
uses out-neighbors, and the directed edge count is the number of arcs from
placed to unplaced vertices; taking maximum or sum gives the four directed
objectives of \cite{wapl14}. Together with the two undirected edge
objectives (Corollary~\ref{cor:edge}(b)) and the two undirected vertex
objectives, these are the eight layout objectives of \cite{wapl14}.

\begin{proposition}[Vertex sum]\label{prop:vsum}
In YES the vertex sum of \(H\) is at most \(2C_K\Lambda(rN+m)\). In NO every
ordering has vertex sum at least \(rNT/10\).
\end{proposition}

\begin{proof}
YES: the ordering of Lemma~\ref{lem:sepYES}(b). NO: fix an ordering. Suppose
first that some edge vertex \(e\) incident to \(u\) appears after at least
\(r/2\) copies of \(u\). During the placements of the copies of \(u\) ranked
\(r/4+1,\dots,r/2\), the first \(r/4\) copies are live (their neighbor
\(e\) is unplaced). This contributes at least \(r^2/16\ge r^2/32\), and
\(r^2/32>rNm/10\ge rNT/10\) by the choice of \(r\). Otherwise every original
vertex has at least \(r/2\) \emph{good} copies, placed after all its incident
edge vertices; let \(g\ge rN/2\) be their number. For each of the at least
\(g/3\ge rN/10\) prefixes containing between \(g/3\) and \(2g/3\) good copies,
there are at least \(rN/6-1\) good copies on each side, hence at least
\(N/6-1/r\ge N/8\) original vertices with a good copy on each side, and every
original vertex has a good copy on some side. Partition \(V(G)=A'\sqcup B'\)
with \(|A'|,|B'|\ge N/8\), \(A'\) having placed good copies and \(B'\)
unplaced ones (assign the vertices with good copies on both sides suitably).
An edge vertex of an edge between \(A'\) and \(B'\) precedes the good copies
of its endpoint in \(A'\), so it is placed, and it has an unplaced neighbor in
\(B'\); so it is live. There are at least \(T\) such edges.
\end{proof}

\begin{proposition}[Interval completion]\label{prop:interval}
For every graph, the minimum vertex sum over orderings equals the minimum
number of edges of an interval supergraph. Consequently, if \(N\ge640C'/B\)
(which Theorem~\ref{th:family} allows), the minimum number of \emph{added}
edges of an interval supergraph of \(H\) is at most \(2C_K\Lambda(rN+m)\) in
YES and at least \(rNT/20\) in NO.
\end{proposition}

\begin{proof}
Given an ordering, give each vertex the interval from its position to the
position of its last later neighbor (or a point). The interval graph
contains the input graph, and the pairs \((u,w)\) with \(u\) earlier and
\(u\)'s interval containing the position of \(w\) are counted exactly by
\(\sum_i\live(i-1)\) at \(w\)'s position \(i\). Conversely, order an interval
supergraph by left endpoints. If \(u\) precedes the vertex \(w\) at
position \(i\) and is live at \(i-1\), with an input neighbor \(v\) at
position \(\ge i\), then \(u\)'s interval meets \(v\)'s, which starts no
earlier than \(w\)'s, so it contains the left endpoint of \(w\); thus
\(uw\) is an edge of the supergraph. Hence the supergraph has at least the
vertex sum of this order many edges. For \(H\), \(|E(H)|=2rm\le2rC'\Lambda
\le rNT/20\) when \(N\ge640C'/B\); subtract from Proposition~\ref{prop:vsum}.
\end{proof}

\begin{proposition}[Directed layouts]\label{prop:directed}
Let \(D_r\) be the DAG on \(V(H_r(G))\) with arcs from each edge vertex to
the copies of its endpoints, for any \(r\ge1\). In YES some topological order
has at every prefix at most \(2C_K\Lambda\) live vertices and at most
\(4rC_K\Lambda\) crossing arcs. In NO, in every topological order at least
\(rN/4\) prefixes have at least \(T\) live vertices and at least \(rT/8\)
crossing arcs.
\end{proposition}

\begin{proof}
YES: the order of Lemma~\ref{lem:sepYES}(b) is topological; each live edge
vertex has at most \(2r\) arcs. NO: consider a prefix containing between
\(rN/3\) and \(2rN/3\) copies; there are at least \(rN/3\) such prefixes. At
least \(N/3\) originals have a placed copy. If \(x\) originals have at least
\(r/8\) unplaced copies, then \(rN/3\le rx+(r/8)(N-x)\), so \(x\ge5N/21\).
Every original lies in one of these two sets (one with no placed copy has
\(r\) unplaced copies), so we may partition \(V(G)=A'\sqcup B'\) with
\(|A'|,|B'|\ge N/8\), every vertex of \(A'\) having a placed copy and every
vertex of \(B'\) at least \(r/8\) unplaced copies. An edge vertex of an edge
between \(A'\) and \(B'\) precedes the placed copy, so it is placed, and it
has at least \(r/8\) arcs to unplaced copies. There are at least \(T\) such
edges.
\end{proof}

\paragraph{Pebbling.}
A one-shot black pebbling of a DAG places a pebble on a vertex only when all
its predecessors carry pebbles, removes pebbles at any time, and pebbles
every vertex exactly once; its cost is the maximum number of simultaneous
pebbles. In one-shot black-white pebbling (complete visiting convention of
\cite{wapl14}) white pebbles may be placed at any time and removed only when
all predecessors carry pebbles, every vertex receives a pebble exactly once,
and every sink is checked, either at its black placement or at the removal
of its white pebble.

\begin{proposition}[Pebbling]\label{prop:pebble}
\begin{enumerate}[label=\textup{(\alph*)}]
\item The one-shot black pebbling cost of a DAG lies between the minimum
over topological orders of the maximum directed live count and that value
plus one. Hence it is at most \(2C_K\Lambda+1\) for \(D_r\) in YES and at
least \(T\) in NO.
\item The same gap holds on DAGs with a single sink and maximum indegree two.
\item For the DAG \(P_H\) with a source \(s_v\) for each vertex \(v\) of \(H\)
and a sink \(t_e\) for each edge \(e\), with the two endpoint sources as
predecessors, the one-shot black-white pebbling cost lies between
\(\pw(H)+1\) and \(\pw(H)+2\).
\end{enumerate}
\end{proposition}

\begin{proof}
(a) The placement order is topological, and a placed vertex with an
unpebbled successor must keep its pebble; so at the placement of the
\(i\)-th vertex all vertices live at \(i-1\) carry pebbles. Conversely,
pebbling in a topological order and removing each pebble after the last
successor costs at most the live maximum plus one. Apply
Proposition~\ref{prop:directed}.

(b) Use the unreplicated DAG \(D_1\). The YES order still has live maximum
at most \(2C_K\Lambda\). In NO, in a topological order of \(D_1\) stop
immediately after half of the original vertices have been placed: the
placed and unplaced originals each number at least \(N/8\), every edge vertex
of an edge between them is placed (it precedes its endpoints) and has an
unplaced successor, so at least \(T\) vertices are live. The sink-merging
tree and indegree-reduction pyramids of \cite[Lemmas~7.1--7.2]{wapl14}
produce a single-sink DAG of maximum indegree two, preserve the lower bound,
and increase the upper bound by at most \(s+1+\max\{2,d\}\), where \(s\le N\)
is the number of sinks and \(d\le N-1\) the maximum indegree of \(D_1\).
Since \(N\le\Lambda\), the YES bound remains \(O(C_K\Lambda)\).

(c) In a one-shot strategy each source carries a pebble during one time
interval. Each sink \(t_{uv}\) is checked at a time when \(s_u\) and \(s_v\)
both carry pebbles. The sets of pebbled sources over time therefore form a
path decomposition of \(H\) (isolated vertices may get singleton bags), with
bags of size at most the cost. Conversely, from a path decomposition of width
\(w\), introduce and remove sources as their bag intervals begin and end
(removing before introducing at each transition), and pebble each sink
\(t_{uv}\) with a black pebble, then remove it, at a bag containing \(u,v\);
this uses at most \(w+2\) pebbles.
\end{proof}

\subsection{Proof of Theorem~\ref{th:main}(A)}

Fix \(c>1\). Take \(K=2\) in Theorem~\ref{th:family} (or \(K\) divisible by
\(k\) or by the denominator of \(\delta\), as in Corollary~\ref{cor:edge}),
so \(C_K\) is fixed, and then a fixed \(B\) with \(B/16>c(C_K+2)\), and
\(n_1\) as required in Proposition~\ref{prop:interval} and for \(N\) large.
With \(s=C_K\Lambda+1\), Lemmas~\ref{lem:sepNO} and~\ref{lem:sepYES} give
\(\td(H)\le s\) in YES and \(\tw(H)\ge B\Lambda/16-1>cs\) in NO; since
\(\tw\le\pw\le\td-1\), this gives every-factor hardness for treewidth,
pathwidth and treedepth. Lemma~\ref{lem:sepYES}(d) and
Proposition~\ref{prop:fillin} give YES \(\le C_K\Lambda^2/2\) versus NO
\(\ge B^2\Lambda^2/2048\) for fill-in; Lemmas~\ref{lem:sepNO}
and~\ref{lem:sepYES}(a) give the balanced vertex separator. The comparison
inequalities, Lemma~\ref{lem:linewidth} and
Proposition~\ref{prop:treealpha} transfer the treewidth gap, up to fixed
constant factors and an additive one that is absorbed by \(\Lambda\), to
branch-width, maximum-matching width, and to clique-, NLC-, rank-,
Boolean-, mim-width and tree-independence number of \(L(H)\).
Propositions~\ref{prop:vsum}--\ref{prop:pebble} and Lemma~\ref{lem:sepYES}(b)
give the vertex layouts, interval completion, the directed layouts and
pebbling, and Corollary~\ref{cor:edge} gives the edge layouts, carving-width
and the partitioning problems. In each case the NO/YES ratio is a fixed
positive multiple of \(B/C_K\) (or of \(B^2/C_K\) for fill-in), so a large
fixed \(B\) gives factor \(c\). For interval completion, a
\(c\)-approximation of the added edges also yields one of the total edges,
by adding back \(|E(H)|\). \qed

\section{Gap small-set expansion}\label{sec:sse}

For a symmetric stochastic \(m\times m\) matrix \(P\) and
\(f\in[0,1]^m\setminus\{0\}\) of \emph{mass} \(z=\frac1m\sum_af_a\), put
\[
\Phi_P(f)=\frac{\sum_{a,b}P_{ab}f_a(1-f_b)}{\sum_af_a}=1-\frac{f^{\top}Pf}{mz}.
\]
For a \(d\)-regular graph with random-walk matrix \(P\), \(\Phi_P(\ind_S)=
\Phi_G(S)\). For a nonnegative symmetric matrix \(Q\) on a set \(V\) and
\(x\in[0,1]^V\), write \(Q[x]=\sum_{i,j}Q_{ij}x_i(1-x_j)\).

\begin{theorem}[Conductance on a size window]\label{th:window}
Fix rationals \(\delta>0\) and \(M\ge1\) with \([\delta/M,M\delta]\subseteq
(0,1/2]\), and \(R>1\). There are fixed \(0<a<b<1\) with \(b/a>R\), a fixed
integer \(d\), and a deterministic polynomial-time reduction from 3SAT to
simple bipartite \(d\)-regular graphs on \(N\) vertices such that: in YES,
some set of exactly \(\delta N\) vertices has conductance at most \(a\); in
NO, every set \(S\) with \(\delta N/M\le|S|\le M\delta N\) has conductance at
least \(b\).
\end{theorem}

\begin{theorem}[Fixed-degree Gap-SSE]\label{th:gapsse}
For every fixed rational \(0<\delta\le1/4\) there are a fixed
\(\eta=\eta(\delta)<1/2\) and a fixed integer \(d\) such that, for each fixed
integer \(j\ge1\), there is a deterministic polynomial-time reduction from
3SAT to simple bipartite \(d\)-regular graphs on \(N\) vertices with: in
YES, some set of exactly \((\delta/j)N\) vertices has conductance at most
\(\eta\); in NO, every nonempty set of at most \((\delta/j)N\) vertices has
conductance at least \(1-\eta\).
\end{theorem}

\subsection{A positive semidefinite chain from the balanced system}

\begin{lemma}\label{lem:chain}
Fix \(\delta,M,R\) as in Theorem~\ref{th:window}, an integer \(K\ge2\) with
\(K\delta\in\Z\), a rational \(B\ge64(R+1)A_K/\delta\), let \(W_1=W_1(K,B)\)
be the bound of Lemma~\ref{lem:balance}, and fix a rational
\(H\ge\max\{2B,32MW_1/\delta\}\). Put \(a_0=A_K/(2H\delta)\) and
\(b_0=B/(16H)\). There is a deterministic polynomial-time reduction from
3SAT to a symmetric, positive semidefinite, stochastic rational matrix \(P\)
of size \(m\times m\) such that: in YES there is \(S'\subseteq[m]\) with
\(|S'|=\delta m\) and \(\Phi_P(\ind_{S'})\le a_0\); in NO,
\(\Phi_P(f)\ge b_0\) for every \(f\in[0,1]^m\) of mass in
\([\delta/M,M\delta]\). Moreover \(b_0/a_0=B\delta/(8A_K)\ge8(R+1)\).
\end{lemma}

\begin{proof}
Let \((V,c,\mu)\) be the output of Lemma~\ref{lem:balance}, \(W=W(c)\le
W_1\), \(h=|V|\), \(d_i=\sum_jc_{ij}\), and
\(\nu_i=(1-2W/H)\mu_i+d_i/H\), a probability vector with
\(\|\nu-\mu\|_{\rm TV}\le2W/H\). Let \(L\) be the symmetric matrix with
\(L_{ij}=c_{ij}/H\) for \(i\ne j\) and \(L_{ii}=(1-2W/H)\mu_i\); its row sums
are \(\nu_i\). Put \(\widehat L=(L+\operatorname{diag}\nu)/2\). It is positive
semidefinite since \(v^{\top}(\operatorname{diag}\nu+L)v=\frac12\sum_{i,j}
L_{ij}(v_i+v_j)^2\ge0\).

\emph{Fractional cuts.} For \(x\in[0,1]^V\) and \(p=\sum_i\mu_ix_i\), let
\(R\subseteq V\) contain \(i\) independently with probability \(x_i\). The
off-diagonal part of \(\widehat L[x]\) equals \(\frac1{2H}\E\,c(\partial
R)\ge\frac B{2H}\bigl[p(1-p)-\sum_i\mu_i^2x_i(1-x_i)\bigr]\) in NO. Since
\(2W/H\le1/2\), the diagonal entry satisfies \(\widehat
L_{ii}\ge\mu_i/4\ge B\mu_i^2/(2H)\) (as \(H\ge2B\)). Hence
\(\widehat L[x]\ge\frac B{2H}p(1-p)\) in NO.

\emph{Copies.} Take \(J\ge16Mh/\delta\), let \(k_i=K\max\{1,\lceil
J\nu_i\rceil\}\), \(m=\sum_ik_i\) and \(\omega=KJ/m\); then \(KJ\le m\le
K(J+h)\) and \(1-\omega\le h/J\le\delta/(16M)\). On the \(m\) copies let
\(Z_{ab}=\omega\widehat L_{ij}/(k_ik_j)\) for \(a\) a copy of \(i\) and \(b\)
a copy of \(j\) (including \(a=b\)). The row sum of \(Z\) at a copy of \(i\)
is \(\omega\nu_i/k_i\le1/m\); add the nonnegative amount
\(1/m-\omega\nu_i/k_i\) on the diagonal, and let \(P\) be \(m\) times the
result. Then \(P\) is symmetric and stochastic, and positive semidefinite
because \(y^{\top}Zy=\omega\,\bar y^{\top}\widehat L\,\bar y\) with
\(\bar y_i=\sum_{a\text{ copy of }i}y_a/k_i\).

\emph{NO.} Let \(f\in[0,1]^m\) have mass \(z\in[\delta/M,M\delta]\), and let
\(x_i\) be the average of \(f\) over the copies of \(i\). Then
\(\sum_{a,b}(P/m)_{ab}f_a(1-f_b)\ge Z[f]=\omega\widehat L[x]\). The class
counting law \(k_i/m\) equals \(\omega\nu\) plus a nonnegative measure of mass
\(1-\omega\), so \(|z-\nu\cdot x|\le1-\omega\), and
\(|z-p|\le(1-\omega)+2W/H\le\delta/(8M)\). On the window this gives
\(p\ge7z/8\) and \(p\le9/16\), so \(p(1-p)\ge z/4\). Therefore
\(\Phi_P(f)\ge\omega\frac B{2H}\cdot\frac z4\cdot\frac1z\ge\frac B{16H}=b_0\),
using \(\omega\ge1/2\).

\emph{YES.} The values \(\nu_i\), hence the counts \(k_i\), are invariant
under \(\mathfrak S_K\), so by Lemma~\ref{lem:balance}(ii) each color class
of keys carries exactly \(m/K\) copies, including copies of zero-mass keys.
Let \(S'\) be the set of copies of the keys in \(K\delta\) color classes;
\(|S'|=\delta m\). The diagonal additions do not contribute to
\(\Phi_P(\ind_{S'})\), and the crossing part is \(\omega\sum c_{ij}/(2H)\)
over pairs of keys of different colors, at most \(A_K/(2H)\). Hence
\(\Phi_P(\ind_{S'})\le A_K/(2H\delta)=a_0\).
\end{proof}

\subsection{Degree reduction}

\begin{lemma}\label{lem:degree}
For every rational \(\eps\in(0,1)\) there is an integer \(D=D(\eps)\) with
the following property. Let \(q\ge D^2\) be a prime and \(L_0\) a square
integer, and let \(P\) be a symmetric stochastic rational \(m\times m\)
matrix. One can compute in time polynomial in \(m,L_0,q\) a simple bipartite
\((D^2q)\)-regular graph \(\Gamma\) whose shores \(V_L,V_R\) have
\(n_s=mL_0q^2\) vertices each, partitioned into \(m\) classes of \(L_0q^2\)
vertices, such that for every \(S\subseteq V(\Gamma)\), with \(\alpha_i\)
and \(\beta_i\) the densities of \(S\) in the \(i\)-th left and right classes
and \(z=|S|/(2n_s)\),
\[
\Bigl|\frac{e(S)}{D^2q\,n_s}-\frac{\alpha^{\top}P\beta}{m}\Bigr|\le
\xi z,\qquad \xi:=q^{-1/2}+2\eps+\frac{2m}{L_0}.
\]
Consequently, for \(S\ne\varnothing\): if \(P\) is positive semidefinite then
\(\Phi_\Gamma(S)\ge\Phi_P(\frac{\alpha+\beta}2)-\xi\); and if
\(\alpha=\beta=\ind_{S'}\) then \(\Phi_\Gamma(S)\le\Phi_P(\ind_{S'})+\xi\).
\end{lemma}

\begin{proof}
\emph{Rounding.} Floor the entries of \(L_0P\) and fill the row and column
deficits greedily, obtaining nonnegative integers \(n_{ij}\) with all row and
column sums \(L_0\); with \(P'=(n_{ij}/L_0)\), every row and column of
\(P'-P\) has \(\ell_1\)-norm at most \(2m/L_0\), hence
\(\|P'-P\|\le2m/L_0\).

\emph{Ports.} Give each left and right class \(L_0\) ports, and fix a
bijection between left and right ports sending exactly \(n_{ij}\) ports of
left class \(i\) to right class \(j\); let \(U\) be its permutation matrix and
\(Z\) the averaging projection onto functions constant on port classes. A
fixed power of the lazy Gabber--Galil expander \cite{gg81} on each square
port class gives a \(D\)-regular multigraph (loops allowed) whose transition
matrix \(Q\) satisfies \(\|Q-Z\|\le\eps\), with \(D\) depending only on
\(\eps\). Then \(\|QUQ-ZUZ\|\le\|(Q-Z)UQ\|+\|ZU(Q-Z)\|\le2\eps\), and \(QUQ\)
is uniform over \(D^2\) walks, so \(v_{ab}=D^2(QUQ)_{ab}\) is an integer
multiplicity at most \(D^2\) with row and column sums \(D^2\).

\emph{Incidence replacement.} Replace each port by a copy of \(\F_q^2\) and,
for each pair of ports \((a,b)\) with \(v_{ab}>0\), choose \(v_{ab}\le q\)
distinct values \(T_{ab}\subseteq\F_q\) and join the two copies by
\(A_{T_{ab}}\). By Lemma~\ref{lem:incidence}(i) the graph is simple,
bipartite and \((D^2q)\)-regular.

\emph{Estimate.} Let \(s_a,t_b\) be the densities of \(S\) at left port \(a\)
and right port \(b\). By Lemma~\ref{lem:incidence}(iii) and
\(\sqrt{s_at_b}\le(s_a+t_b)/2\),
\(\bigl|e(S)-q^3\sum_{a,b}v_{ab}s_at_b\bigr|\le q^{5/2}\sum_{a,b}v_{ab}
(s_a+t_b)/2=q^{5/2}D^2(\sum_as_a+\sum_bt_b)/2\). Dividing by
\(D^2q\,n_s=D^2q^3mL_0\) and noting \(z=(\sum_as_a+\sum_bt_b)/(2mL_0)\),
\[
\Bigl|\frac{e(S)}{D^2qn_s}-\frac{s^{\top}QUQt}{mL_0}\Bigr|\le zq^{-1/2}.
\]
Next \(|s^{\top}(QUQ-ZUZ)t|/(mL_0)\le2\eps\|s\|\|t\|/(mL_0)\le
\eps(\sum s_a+\sum t_b)/(mL_0)=2\eps z\). Since \(Zs\) and \(Zt\) take the
values \(\alpha_i\) and \(\beta_j\) on the port classes,
\(s^{\top}ZUZt=\sum_{i,j}n_{ij}\alpha_i\beta_j=L_0\alpha^{\top}P'\beta\).
Finally \(|\alpha^{\top}(P'-P)\beta|/m\le\frac{2m}{L_0}\cdot\frac{\|\alpha\|
\|\beta\|}m\le\frac{2m}{L_0}\cdot\frac{\sum\alpha_i+\sum\beta_i}{2m}=
\frac{2m}{L_0}z\).

\emph{Consequences.} \(\Phi_\Gamma(S)=1-2e(S)/(D^2q|S|)=1-e(S)/(D^2qn_sz)\),
and \(f=(\alpha+\beta)/2\) has \(\sum_if_i=mz\). With \(g=(\alpha-\beta)/2\)
and \(P\) symmetric, \(\alpha^{\top}P\beta=f^{\top}Pf-g^{\top}Pg\le
f^{\top}Pf\) when \(P\) is positive semidefinite, which gives the first
consequence; the second follows from the lower estimate with
\(\alpha=\beta=\ind_{S'}\).
\end{proof}

\subsection{Proofs of Theorems~\ref{th:window} and~\ref{th:gapsse}}

\begin{proof}[Proof of Theorem~\ref{th:window}]
Fix \(K\) with \(K\delta\in\Z\), then \(B\) and \(H\) as in
Lemma~\ref{lem:chain}. Fix \(\eps\) and a prime \(q\ge D(\eps)^2\) with
\(q^{-1/2}+2\eps\le a_0/4\), and take a square \(L_0\), polynomial in \(m\),
with \(2m/L_0\le a_0/4\); so \(\xi\le a_0/2\). Apply Lemma~\ref{lem:degree}
to the matrix \(P\) of Lemma~\ref{lem:chain}; \(N=2n_s\) and \(d=D^2q\) is
fixed. In YES, the set \(S\) consisting of both shores of the classes in
\(S'\) has exactly \(\delta N\) vertices and \(\Phi_\Gamma(S)\le
\frac32a_0\le a:=2a_0\). In NO, a set with \(\delta N/M\le|S|\le M\delta N\)
has \(z\) in the window and \(\Phi_\Gamma(S)\ge b_0-a_0/2\ge b:=b_0/2\),
since \(a_0\le b_0/8\). Finally \(b/a=b_0/(4a_0)\ge2(R+1)>R\).
\end{proof}

\begin{proof}[Proof of Theorem~\ref{th:gapsse}]
Take \(M=2\) and \(R=2\) in Lemma~\ref{lem:chain}, and \(H\) so large that
\(b_0\le\delta/100\) (this lowers \(a_0\) and \(b_0\) proportionally). Let
\(\kappa=(a_0+b_0)/2\) and
\[
\lambda=\frac{1/2-\kappa}{1-\delta-\kappa}\in(0.4,1),\qquad\text{so that}\qquad
(1-\lambda)\kappa+\lambda(1-\delta)=\tfrac12,
\]
and \(P_\lambda=(1-\lambda)P+\lambda J/m\), again symmetric, stochastic and
positive semidefinite. For \(f\) of mass \(z\),
\(\Phi_{J/m}(f)=1-z\), so \(\Phi_{P_\lambda}(f)=(1-\lambda)\Phi_P(f)+
\lambda(1-z)\). Put \(g=(1-\lambda)(b_0-a_0)/2>0\). In YES,
\(\Phi_{P_\lambda}(\ind_{S'})\le(1-\lambda)a_0+\lambda(1-\delta)=1/2-g\). In
NO, for \(z\in[\delta/2,\delta]\), \(\Phi_{P_\lambda}(f)\ge(1-\lambda)b_0+
\lambda(1-\delta)=1/2+g\); for \(0<z<\delta/2\), using only
\(\Phi_P(f)\ge0\), \(\Phi_{P_\lambda}(f)\ge\lambda(1-\delta)+\lambda\delta/2
\ge\lambda(1-\delta)+(1-\lambda)b_0=1/2+g\), because
\(\lambda\delta/2\ge0.2\delta\ge(1-\lambda)b_0\). Apply
Lemma~\ref{lem:degree} to \(P_\lambda\) with \(\xi<g/2\) and put
\(\eta=1/2-g/2\): YES gives a set of exactly \(\delta N\) vertices with
conductance at most \(\eta\), and NO gives conductance at least \(1-\eta\)
for every nonempty set of at most \(\delta N\) vertices. For density
\(\delta/j\), output \(j\) disjoint copies of this graph: the YES set of one
copy has \((\delta/j)\) of all vertices; in NO a set of at most
\(\delta N\) vertices meets each copy in at most \(\delta N\) vertices, and
its conductance is the size-weighted mean of the conductances of its
nonempty parts. The degree and \(\eta\) do not depend on \(j\).
\end{proof}

This proves Theorem~\ref{th:main}(B). In both theorems the YES set takes
the same number of vertices from each shore.

\section{Bicliques}\label{sec:biclique}

For a bipartite graph with shores \(V_L,V_R\), a \emph{biclique} is a pair
\(X\subseteq V_L\), \(Y\subseteq V_R\) with all pairs adjacent; its shore
densities are \(|X|/|V_L|\) and \(|Y|/|V_R|\). Maximum Balanced Biclique
asks for the largest \(k\) with a biclique of shores of size \(k\); Maximum
Edge Biclique asks for the maximum of \(|X||Y|\).

\begin{lemma}[Base gap]\label{lem:basegap}
For every fixed \(\xi>0\) there are fixed rationals \(b<a<1/2\) with
\(a\ge1/2-\xi\) and a deterministic polynomial-time reduction from 3SAT to
bipartite graphs with equal shores such that in YES some biclique has two
equal shores of density at least \(a\), and in NO every biclique has shore
densities \(\alpha,\beta\) with \(\sqrt{\alpha\beta}\le b\).
\end{lemma}

\begin{proof}
Apply Lemma~\ref{lem:balance} with \(K=2\) and write \(A=A_2\). Fix \(B\ge
64A\), let \(W_*=W_1(2,B)\), and fix \(H\ge100\max\{W_*,B,A\}\) with
\(A/H\le\xi\). Let \((V,c,\mu)\) be the output and \(W=W(c)\le W_*\).
Form a \emph{conflict graph}: its vertices are the keys \(i\in V\) and one
port \(e\) for each pair \(e=\{i,j\}\) with \(c_e>0\), the port being
adjacent to \(i\) and \(j\). Give keys and ports the weights
\[
w_i=(1-3W/H)\mu_i+d_i/H,\qquad w_e=c_e/H,
\]
which are nonnegative and sum to \((1-3W/H)+2W/H+W/H=1\). Let \(s\) be the
number of keys and ports, fix \(\eps=\min\{A/H,1/100\}\) and take
\(J\ge s/\eps\). Replace each key or port \(u\) by a clique of
\(Jw_u+e_u\) copies, \(0\le e_u\le1\) (namely \(\max\{1,\lceil Jw_u\rceil\}\)),
and join copies of adjacent classes completely. The counting distribution on
classes is a mixture of \(w\) and the normalized rounding excess, of excess
fraction at most \(s/J\le\eps\), so it is within \(\eps\) of \(w\) in total
variation. Let \(m\) be the number of copies. The output \(\Gamma\) has two
shores, each a copy of this set of \(m\) copies, and \(x\in V_L\),
\(y\in V_R\) are adjacent iff they are distinct and not adjacent in the
conflict graph. Thus equality counts as a conflict.

\emph{YES.} Let \(\chi\) be the coloring of Lemma~\ref{lem:balance}(ii).
Omit the ports of pairs \(\{i,j\}\) with \(\chi(i)\ne\chi(j)\), and let
\(X\) (left) consist of the copies of keys of color \(1\) and of ports with
both ends of color \(1\), and \(Y\) (right) likewise for color \(2\). Keys
are adjacent in the conflict graph only to ports, and the retained ports are
adjacent only to keys of their own color, so \((X,Y)\) is a biclique. The
color swap preserves degrees, capacities, weights and hence the copy counts,
and it maps the left set onto the right set; so if \(h\) of the \(m\) copies
are omitted, both shores have exactly \((m-h)/2\) vertices. The omitted
weight is at most \(A/H\), so \(h/m\le A/H+\eps\le2A/H\), and both shore
densities are at least \(a:=1/2-A/H\ge1/2-\xi\).

\emph{NO.} Let \((X,Y)\) be a biclique of densities \(\alpha,\beta\). The
underlying copies of \(X\) and \(Y\) are disjoint, and since each class is a
clique of conflicts, no class is represented on both shores; so
\(\alpha+\beta\le1\). If \(\min\{\alpha,\beta\}<1/3\), then
\(\alpha\beta<\frac13\cdot\frac23=\frac29<b^2\) for the \(b\) below (as
\(B/H\le1/100\)). Suppose \(\alpha,\beta\ge1/3\), and let \(X_0,Y_0\) be the
sets of keys with a copy in \(X\), respectively \(Y\), and \(U\) the keys
represented on neither side. Since the ports and degree terms weigh \(3W/H\)
in total, \(\mu(X_0),\mu(Y_0)\ge1/3-3W/H-\eps\ge1/4\), so
\(\mu(X_0)\in[1/4,3/4]\). Every port between \(X_0\) and \(Y_0\) conflicts
with a copy in \(X\) and one in \(Y\), so it is omitted; every key in \(U\)
is omitted. These classes are distinct, and their weight is at least
\[
\frac{c(X_0,Y_0)+\sum_{u\in U}d_u}{H}\ \ge\ \frac{c(\partial X_0)}{H}
\ \ge\ \frac{B\mu(X_0)(1-\mu(X_0))}H\ \ge\ \frac{3B}{16H}.
\]
Passing to counts loses at most \(\eps\le A/H\le B/(64H)\), so at least a
fraction \(B/(8H)\) of the copies is omitted, and
\(\sqrt{\alpha\beta}\le(\alpha+\beta)/2\le\frac12-\frac B{16H}=:b\). Since
\(B/16>A\), \(b<a\).
\end{proof}

\begin{remark}
The same omitted-weight argument, with \(1/3\) replaced by a fixed balance
threshold below \(1/2\), shows that the minimum number of vertices whose
deletion from the blown-up conflict graph leaves two anticomplete sides,
each of at least that fraction of the vertices, is NP-hard to approximate
within every fixed factor (taking \(B/A\) large).
\end{remark}

\begin{theorem}[Polynomial-factor hardness]\label{th:biclique}
Assume Theorem~U. There is a fixed \(\gamma>0\) such that Maximum Balanced
Biclique and Maximum Edge Biclique are NP-hard to approximate within
\(N^\gamma\) on bipartite graphs with \(N\) vertices.
\end{theorem}

\begin{proof}
Take Lemma~\ref{lem:basegap} with any \(\xi<1/2\), and fix \(a_1\) with
\(b<a_1<a\). Pad each shore of size \(m\) to \(n=\lceil\sqrt m\rceil^2\)
with isolated vertices; a fixed lower bound on \(J\) ensures \(am/n\ge
a_1\), and NO densities can only decrease (isolated vertices lie in no
biclique with a nonempty other shore). Identify each shore with
\([n]=\Z_{\sqrt n}^2\).

Fix \(\lambda\in(0,1]\) with \(\Theta:=(a_1-\lambda)/(b+\sqrt{3\lambda})>1\),
and take a fixed power of the lazy Gabber--Galil walk \cite{gg81} on \([n]\):
a symmetric positive semidefinite transition matrix \(T\) of constant labeled
degree \(D\) with \(\|T-J/n\|\le\lambda\). Let \(k=\lceil\log_2n\rceil\). A
\emph{record} is a starting symbol together with \(k\) edge labels; it
determines a walk \((v_0,\dots,v_k)\), and the uniform measure on the
\(nD^k\) records is the stationary walk law (parallel transitions and loops
counted with multiplicity). The output graph has a shore of \(nD^k\) records
on each side, and two records are adjacent iff their walks are adjacent in
the padded base graph at every position. There are \(N=2nD^k\le2(2D)^k\)
vertices, enumerable in polynomial time.

\emph{YES.} For a base shore set \(S\) of density \(\alpha\ge a_1\), let
\(P_S\) be the diagonal projection onto \(S\). The fraction of records
staying in \(S\) is \(n^{-1}\ind_S^{\top}(P_STP_S)^k\ind_S\). The compressed
operator is positive semidefinite and its Rayleigh quotient at \(\ind_S\) is
\(\ind_S^{\top}T\ind_S/|S|\ge\alpha-\lambda\); Jensen's inequality on the
spectral measure of \(\ind_S\) gives a fraction at least
\(\alpha(\alpha-\lambda)^k\). Records staying in the two shores of the YES
biclique form a biclique, so after trimming the larger side there is a
balanced biclique with shore density at least \(y=a_1(a_1-\lambda)^k\).

\emph{NO.} Let \((\mathcal X,\mathcal Y)\) be a biclique of records and
\(S_i,U_i\) the sets of position-\(i\) symbols on the two sides. Every pair in
\(S_i\times U_i\) is a base edge (choose the witnessing records
independently), so the densities satisfy \(\alpha_i\beta_i\le b^2\). Since
\(\|P_{S_i}TP_{S_{i+1}}\|\le\|P_{S_i}(J/n)P_{S_{i+1}}\|+\lambda=
\sqrt{\alpha_i\alpha_{i+1}}+\lambda\), the density of \(\mathcal X\) is at most
\(n^{-1}\ind_{S_0}^{\top}P_{S_0}TP_{S_1}T\cdots TP_{S_k}\ind_{S_k}\le
\sqrt{\alpha_0\alpha_k}\prod_{i<k}(\sqrt{\alpha_i\alpha_{i+1}}+\lambda)\), and
likewise for \(\mathcal Y\). With \(u=\sqrt{\alpha_i\alpha_{i+1}}\) and
\(v=\sqrt{\beta_i\beta_{i+1}}\) we have \(uv\le b^2\) and
\((u+\lambda)(v+\lambda)\le b^2+2\lambda+\lambda^2\le(b+\sqrt{3\lambda})^2\).
The endpoint factors are at most one, so the geometric mean of the two shore
densities is at most \(z=(b+\sqrt{3\lambda})^k\); this covers unbalanced
bicliques. (This estimate also follows from \cite[Proposition~22]{tz24}
applied to \(T\otimes T\); see also \cite[Section~4]{gl04}.)

The ratio is \(y/z=a_1\Theta^k\). For every fixed \(0<\gamma<\log\Theta/
\log(2D)\), \(a_1\Theta^k>N^\gamma\) for all large \(n\) (a fixed lower bound
on \(J\) ensures this). For Maximum Balanced Biclique, YES has value at least
\(nD^ky\) and NO at most \(nD^kz\). For Maximum Edge Biclique, YES has a
biclique with at least \((nD^ky)^2\) edges and every NO biclique has at most
\((nD^kz)^2\); the ratio is squared.
\end{proof}

This proves Theorem~\ref{th:main}(C).

\section{Dense subgraphs and cuts}\label{sec:dense}

This section uses the constructions of Section~\ref{sec:sse}.

\begin{lemma}[Window graphs of large degree]\label{lem:largedeg}
Fix \(\delta,M,R\) as in Theorem~\ref{th:window}, and let \(a,b\) be the
constants of its proof (\(a=2a_0\), \(b=b_0/2\)). For every fixed
\(\theta\ge1\) there is a deterministic polynomial-time reduction from 3SAT
to simple bipartite \(d\)-regular graphs on \(N\) vertices with
\(d\ge\theta\sqrt N\) and \(N\ge\theta d\), \(\delta N\ge\theta^2\), with the
YES and NO properties of Theorem~\ref{th:window} for these \(a,b\).
\end{lemma}

\begin{proof}
Let \(P\) be as in Lemma~\ref{lem:chain}. Take \(s\ge\max\{m,\theta^2,s_0\}\)
for a fixed \(s_0\), an integer \(L_0\) of order \(s^2\) and a prime
\(q\in[sL_0,2sL_0]\). Round \(P\) to \((n_{ij}/L_0)\) as in
Lemma~\ref{lem:degree}, and build \(\Gamma\) directly: each left and right
class is a copy of \(\F_q^2\), and left class \(i\) is joined to right class
\(j\) by \(A_{T_{ij}}\) with \(n_{ij}\le L_0\le q\) distinct values. It is
simple, bipartite and \(d\)-regular with \(d=L_0q\), and \(N=2mq^2\). The
estimate of Lemma~\ref{lem:degree}, with the port step omitted, gives the
error \(\xi=q^{-1/2}+2m/L_0\to0\); for \(s_0\) large, \(\xi\le a_0/2\), and
the proof of Theorem~\ref{th:window} applies. Moreover
\(d/\sqrt N=L_0/\sqrt{2m}\) and \(N/d=2mq/L_0\ge2ms\), both at least
\(\theta\) for \(s_0\) large, and \(\delta N\ge\theta^2\).
\end{proof}

In a \(d\)-regular graph,
\begin{equation}\label{eq:edges}
e(S)=\tfrac12d|S|\bigl(1-\Phi(S)\bigr).
\end{equation}

\begin{theorem}[Densest \(k\)-Subgraph]\label{th:dks}
Let \(a,b\) be as in Theorem~\ref{th:window} for some fixed window. Maximizing
\(e(S)\) over \(|S|=k\) is NP-hard to approximate within any fixed factor
below \((1-a)/(1-b)\) on fixed-degree simple bipartite graphs, also when
\(S\) must contain exactly \(k/2\) vertices of each shore.
\end{theorem}

\begin{proof}
Use the graphs of Theorem~\ref{th:window} with \(k=\delta N\). By
\eqref{eq:edges}, YES has a \(k\)-set with \(e(S)\ge\frac12dk(1-a)\), taking
\(k/2\) vertices from each shore; in NO every \(k\)-set has
\(e(S)\le\frac12dk(1-b)\).
\end{proof}

\begin{theorem}[Densest At-Least-\(k\)-Subgraph]\label{th:dalks}
Let \(M\ge2\) and \(a,b\) be as in Theorem~\ref{th:window}. Maximizing
\(e(S)/|S|\) over \(|S|\ge k\) is NP-hard to approximate within any fixed
factor below \(1+(b-a)/(1+b/(M-1))\) on simple unweighted graphs.
\end{theorem}

\begin{proof}
Use Lemma~\ref{lem:largedeg} with a large fixed \(\theta\), and let
\(k=\delta N\). Put \(t=bM/(M-1)\), let \(r\) be the largest integer with
\(\binom r2\le tdk/2\), and add a disjoint clique on \(r\) vertices; it has
\(tdk/2-O(r)\) edges and \(r=O(\sqrt{dk})\), so \(r/k=O(\sqrt{d/k})\to0\) as
\(\theta\to\infty\) (since \(d/k\le1/(\delta\theta)\)). Set the new size bound
to \(k+r\). In units of \(d/2\), the YES witness together with the clique
has density at least \((1-a+t)\frac{k}{k+r}-O(r/(dk))\). A NO solution
contains \(u\ge k\) old vertices. If \(u\le Mk\), by \eqref{eq:edges} and the
window its density is at most \((u(1-b)+tk)/u\le1-b+t\); if \(u>Mk\),
regularity gives at most \((u+tk)/u\le1+t/M\). The two bounds agree, and
\[
\frac{1-a+t}{1+t/M}=1+\frac{b-a}{1+b/(M-1)} .
\]
A fixed factor strictly below this value absorbs the \(O(r/k)\) losses for
\(\theta\) large.
\end{proof}

\begin{theorem}[Minimum \(k\)-Cut]\label{th:kcut}
Let \(a,b\) be as in Theorem~\ref{th:window} with \(M\ge2\). Minimizing the
number of edges between the parts of a partition into \(k\) nonempty parts
(\(k\) part of the input) is NP-hard to approximate within any fixed factor
below \((1+b)/(1+a)\) on connected simple unweighted graphs.
\end{theorem}

\begin{proof}
Use Lemma~\ref{lem:largedeg} with a large fixed \(\theta\), let
\(k=\delta N\), and ask for \(k+1\) parts. In YES, making each vertex of the
YES set a singleton part and the rest one part costs
\(dk-e(S)\le\frac12(1+a)dk\). In NO, let \(r=\lceil\sqrt N\rceil\) and sort
the parts by size, \(|C_1|\le\dots\le|C_{k+1}|\). Take the longest initial
segment \(C_1,\dots,C_j\) with \(j\le k-r\) parts and at most \(k\) vertices,
and let \(U\) be its union. At least \(r+1\) parts follow \(C_j\), so
\((r+1)|C_j|\le N\) and every selected part has fewer than \(r\) vertices.
If \(j=k-r\), then \(|U|\ge k-r\). Otherwise adding \(C_{j+1}\) would exceed
\(k\) vertices, and \(C_{j+1}\) is followed by at least \(r+1\) parts, so
\(|C_{j+1}|<r\) and \(|U|>k-r\). Thus \(k-r\le|U|\le k\), within the window
for large \(N\). Edges leaving \(U\) and edges inside \(U\) between distinct
selected parts all cross the partition, so the cost is at least
\[
d|U|-e(U)-\sum_{j'\le j}e(C_{j'})\ \ge\ d|U|-\tfrac12d|U|(1-b)-\tfrac12r|U|
=\tfrac12\bigl(d(1+b)-r\bigr)|U|,
\]
using \eqref{eq:edges} for \(U\) and \(e(C)\le|C|r/2\). As \(r/d\le
1/\theta+1/\sqrt N\) and \(r/k\to0\), the ratio tends to \((1+b)/(1+a)\) for
large \(\theta\). To make the graph connected, add at most \(N-1\) edges
joining its components. This adds at most \(N-1\) to the YES cost, a
vanishing fraction of \(dk\) since \(N/(dk)=1/(\delta d)\le1/(\delta\theta)\),
and adding edges cannot decrease the cost of any partition.
\end{proof}

\begin{proof}[Proof of Theorem~\ref{th:main}(D)]
Fix one window (say \(\delta=1/8\), \(M=2\), \(R=2\)) with its constants
\(a<b\), and take \(\eps_0>0\) strictly below the three margins
\((1-a)/(1-b)-1\), \((b-a)/(1+b/(M-1))\) and \((b-a)/(1+a)\) of
Theorems~\ref{th:dks}--\ref{th:kcut}.
\end{proof}

\section{Open problems}\label{sec:open}

\paragraph{Full small-set expansion.}
Theorem~\ref{th:gapsse} has one threshold \(\eta(\delta)\) close to \(1/2\);
the SSE hypothesis \cite{rs10} asks for every small \(\eta\). The
random-walk amplification of Raghavendra and Schramm \cite{rsch14} needs
soundness asymptotically larger than the square root of the completeness,
on an enlarged range of sizes. Our normalization gives
\(b_0^2/a_0=B^2\delta/(128A_KH)\), which diverges only if \(H=o(B^2)\), while
Lemma~\ref{lem:chain} uses \(H\ge2B\) and \(H\) of the order of the total
capacity. A rounding argument based on Trevisan's product-demand Cheeger
approximation \cite{tre13} suggests a barrier: a reduction that keeps a
partition of the YES demand into classes of small mass with bounded crossing
capacity \(A\), and imposes \(c(\partial S)\ge B\mu(S)(1-\mu(S))\) on a fixed
size window, needs total capacity \(\Omega(B^2/A)\) unless
\(\mathrm P=\mathrm{NP}\). (We do not include the proof here.) Full SSE by
this route would therefore need a different completeness promise.

\paragraph{The near-half biclique promise.}
Manurangsi's SSE-based base gap \cite{man18} has YES bicliques of density near
\(1/2\) and NO bicliques of arbitrarily small density. Lemma~\ref{lem:basegap}
gives only fixed \(b<a\), and the product of Section~\ref{sec:biclique}
decreases the YES density. A symmetric ``OR'' product does not help either.
Declare two \(r\)-tuples compatible if every bijection between their
positions pairs some compatible base symbols. Consider a conflict graph on
three cliques \(C,T,D\) of masses \(p\), \(vp\), \(1-(1+v)p\), with \(T\)
adjacent to everything and no conflicts between \(C\) and \(D\). Every base
biclique has one shore of mass at most \(p\) and omits mass at least \(vp\).
In the product, the tuples with some entry in \(C\) and the tuples entirely
in \(D\) form a biclique of densities \(1-(1-p)^r\) and \((1-(1+v)p)^r\),
which tend to \(1/2\) and \(2^{-(1+v)}\) for \(p=(\log2)/r\).

\paragraph{\(2-\eps\) for precedence-constrained scheduling.}
The near-two lower bounds for precedence-constrained scheduling of Svensson
\cite{sve11} and of Bazzi and Norouzi-Fard \cite{bnf15} rely on the near-half
biclique promise or a related availability promise, and remain open from
Theorem~U. Likewise, the \(2-\eps\) bounds for Minimum \(k\)-Cut and Densest
At-Least-\(k\)-Subgraph \cite{man18} need conductance parameters stronger
than those of Theorem~\ref{th:window}.

\section*{Use of AI}
The proof presented in this paper was developed through an interactive
dialogue with GPT-6.2 Astra (OpenAI) and Claude Opus 5.5 (Anthropic).
The author used them to explore multiple possible proof strategies and
directions, and the AI generated the initial draft under the author's
guidance. The author then critically examined the generated proof,
verified each step for correctness and rigor, and revised it where
necessary. The author takes full responsibility for the final content.

\begin{thebibliography}{99}

\bibitem{afwz95}
N.~Alon, U.~Feige, A.~Wigderson, and D.~Zuckerman.
\href{https://www.math.ias.edu/avi/node/790}{Derandomized graph products}.
\emph{Comput. Complexity} 5 (1995), 60--75.

\bibitem{ams11}
C.~Amb\"uhl, M.~Mastrolilli, and O.~Svensson.
\href{https://theory.epfl.ch/osven/Ola%20Svensson_publications/SICOMP11a.pdf}{Inapproximability results for maximum edge biclique, minimum linear arrangement, and sparsest cut}.
\emph{SIAM J. Comput.} 40 (2011), 567--596.

\bibitem{ar06}
K.~Andreev and H.~R\"acke.
\href{https://www.cs.cmu.edu/~harry/pdf/bpartition.pdf}{Balanced graph partitioning}.
\emph{Theory Comput. Syst.} 39 (2006), 929--939.

\bibitem{apw12}
P.~Austrin, T.~Pitassi, and Y.~Wu.
\href{https://arxiv.org/abs/1109.4910}{Inapproximability of treewidth, one-shot pebbling, and related layout problems}.
APPROX 2012; arXiv:1109.4910.

\bibitem{bfk14}
N.~Bansal, U.~Feige, R.~Krauthgamer, K.~Makarychev, V.~Nagarajan, J.~Naor, and R.~Schwartz.
\href{https://doi.org/10.1137/120873996}{Min-max graph partitioning and small set expansion}.
\emph{SIAM J. Comput.} 43 (2014), 872--904.

\bibitem{bnf15}
A.~Bazzi and A.~Norouzi-Fard.
\href{https://arxiv.org/abs/1507.01906}{Towards tight lower bounds for scheduling problems}.
ESA 2015; arXiv:1507.01906.

\bibitem{bccfv10}
A.~Bhaskara, M.~Charikar, E.~Chlamt\'a\v{c}, U.~Feige, and A.~Vijayaraghavan.
\href{https://users.cs.northwestern.edu/~aravindv/DkS.pdf}{Detecting high log-densities: an \(O(n^{1/4})\) approximation for densest \(k\)-subgraph}.
STOC 2010, 201--210.

\bibitem{bhangale16}
A.~Bhangale, R.~Gandhi, M.~T.~Hajiaghayi, R.~Khandekar, and G.~Kortsarz.
\href{https://drops.dagstuhl.de/entities/document/10.4230/LIPIcs.ICALP.2016.6}{Bicovering: covering edges with two small subsets of vertices}.
ICALP 2016, LIPIcs 55, article 6.

\bibitem{bon}
\'E.~Bonnet.
\href{https://arxiv.org/abs/2406.11628}{Treewidth inapproximability and tight ETH lower bound}.
arXiv:2406.11628 (17 June 2024); STOC 2025;
\emph{J. ACM} (2026), \href{https://doi.org/10.1145/3833387}{doi:10.1145/3833387}.

\bibitem{bns25}
\'E.~Bonnet, D.~Neuen, and M.~Soko\l{}owski.
\href{https://drops.dagstuhl.de/entities/document/10.4230/LIPIcs.IPEC.2025.17}{Treedepth inapproximability and exponential ETH lower bound}.
IPEC 2025, LIPIcs 358, article 17; arXiv:2507.13818.

\bibitem{bmpy}
N.~Brettell, A.~Munaro, D.~Paulusma, and S.~Yang.
\href{https://arxiv.org/abs/2308.05817}{Comparing width parameters on graph classes}.
\emph{European J. Combin.} (2025), \href{https://doi.org/10.1016/j.ejc.2025.104163}{doi:10.1016/j.ejc.2025.104163}; arXiv:2308.05817.

\bibitem{cs17}
Y.~Cao and R.~B.~Sandeep.
\href{https://arxiv.org/abs/1606.08141}{Minimum fill-in: inapproximability and almost tight lower bounds}.
SODA 2017; arXiv:1606.08141.

\bibitem{dfgkm22}
C.~Dallard, F.~V.~Fomin, P.~A.~Golovach, T.~Korhonen, and M.~Milani\v{c}.
\href{https://arxiv.org/abs/2207.09993}{Computing tree decompositions with small independence number}.
arXiv:2207.09993.

\bibitem{dms21}
C.~Dallard, M.~Milani\v{c}, and K.~\v{S}torgel.
\href{https://arxiv.org/abs/2111.04543}{Treewidth versus clique number. III. Tree-independence number of graphs with a forbidden structure}.
arXiv:2111.04543.

\bibitem{es81}
B.~Efron and C.~Stein.
The jackknife estimate of variance.
\emph{Ann. Statist.} 9 (1981), 586--596.

\bibitem{fei02}
U.~Feige.
\href{https://doi.org/10.1145/509907.509985}{Relations between average case complexity and approximation complexity}.
STOC 2002, 534--543.

\bibitem{fk04}
U.~Feige and S.~Kogan.
\href{https://www.wisdom.weizmann.ac.il/~feige/TechnicalReports/FeigeKogan.pdf}{Hardness of approximation of the balanced complete bipartite subgraph problem}.
Technical report, Weizmann Institute, May 2004.

\bibitem{gg81}
O.~Gabber and Z.~Galil.
\href{https://doi.org/10.1016/0022-0000(81)90040-4}{Explicit constructions of linear-sized superconcentrators}.
\emph{J. Comput. System Sci.} 22 (1981), 407--420.

\bibitem{gl22}
S.~Ghoshal and E.~Lee.
\href{https://arxiv.org/abs/2201.04617}{A characterization of approximability for biased CSPs}.
arXiv:2201.04617.

\bibitem{gl04}
A.~Goerdt and A.~Lanka.
\href{https://eccc.weizmann.ac.il/report/2004/048/}{An approximation hardness result for bipartite clique}.
ECCC TR04-048, 2004.

\bibitem{gw07}
F.~Gurski and E.~Wanke.
\href{https://www.sciencedirect.com/science/article/pii/S0012365X07000684}{Line graphs of bounded clique-width}.
\emph{Discrete Math.} 307 (2007), 2734--2754.

\bibitem{icalp25}
B.~Bergougnoux, \'E.~Bonnet, and J.~Duron.
\href{https://doi.org/10.4230/LIPIcs.ICALP.2025.25}{Mim-width is paraNP-complete}.
ICALP 2025, LIPIcs 334, article 25.

\bibitem{icalp26}
M.~Dupr\'e la Tour, M.~Lafond, and N.~Ndiaye.
\href{https://doi.org/10.4230/LIPIcs.ICALP.2026.83}{On the hardness of recognizing graphs of small mim-width and its variants}.
ICALP 2026, LIPIcs 374, article 83; arXiv:2512.06186.

\bibitem{jst15}
L.~Jeong, S.~H.~S\ae{}ther, and J.~A.~Telle.
\href{https://drops.dagstuhl.de/entities/document/10.4230/LIPIcs.IPEC.2015.212}{Maximum matching width: new characterizations and a fast algorithm for dominating set}.
IPEC 2015; arXiv:1507.02384.

\bibitem{joh98}
\"O.~Johansson.
Clique-decomposition, NLC-decomposition, and modular decomposition---relationships and results for random graphs.
\emph{Congr. Numer.} 132 (1998), 39--60.

\bibitem{kho06}
S.~Khot.
\href{https://doi.org/10.1137/S0097539705447037}{Ruling out PTAS for graph min-bisection, dense \(k\)-subgraph, and bipartite clique}.
\emph{SIAM J. Comput.} 36 (2006), 1025--1071; FOCS 2004.

\bibitem{kin92}
N.~G.~Kinnersley.
The vertex separation number of a graph equals its path-width.
\emph{Inform. Process. Lett.} 42 (1992), 345--350.

\bibitem{lmt26}
B.~Laekhanukit, P.~Manurangsi, and O.~Trabelsi.
A note on approximability of densest at-least-\(k\)-subgraph.
\emph{Theoret. Comput. Sci.} (2026), article 116209,
\href{https://doi.org/10.1016/j.tcs.2026.116209}{doi:10.1016/j.tcs.2026.116209};
\href{https://arxiv.org/abs/2605.25464}{arXiv:2605.25464}.

\bibitem{man17}
P.~Manurangsi.
\href{https://arxiv.org/abs/1611.05991}{Almost-polynomial ratio ETH-hardness of approximating densest \(k\)-subgraph}.
STOC 2017; revised version \href{https://eccc.weizmann.ac.il/report/2016/195/}{ECCC TR16-195} (October 2026).

\bibitem{man17icalp}
P.~Manurangsi.
\href{https://doi.org/10.4230/LIPIcs.ICALP.2017.79}{Inapproximability of maximum edge biclique, maximum balanced biclique and minimum \(k\)-cut from the small set expansion hypothesis}.
ICALP 2017, LIPIcs 80, article 79.

\bibitem{man18}
P.~Manurangsi.
\href{https://www.mdpi.com/1999-4893/11/1/10}{Inapproximability of maximum biclique problems, minimum \(k\)-cut and densest at-least-\(k\)-subgraph from the small set expansion hypothesis}.
\emph{Algorithms} 11 (2018), article 10; arXiv:1705.03581.

\bibitem{man26}
P.~Manurangsi.
Improved approximation algorithm for maximum balanced biclique.
APPROX/RANDOM 2026, LIPIcs 392, article 1,
\href{https://doi.org/10.4230/LIPIcs.APPROX/RANDOM.2026.1}{doi:10.4230/LIPIcs.APPROX/RANDOM.2026.1}; arXiv:2604.27141.

\bibitem{oai26}
OpenAI.
Constant-factor hardness of uniform sparsest cut.
Preprint, 24 September 2026. \url{https://github.com/openai/math}.

\bibitem{oailean}
OpenAI.
Lean formalization scope for the uniform sparsest cut release.
\url{https://github.com/openai/math/blob/main/lean/docs/117.md}.

\bibitem{pee03}
R.~Peeters.
\href{https://doi.org/10.1016/S0166-218X(03)00333-0}{The maximum edge biclique problem is NP-complete}.
\emph{Discrete Appl. Math.} 131 (2003), 651--654.

\bibitem{rsch14}
P.~Raghavendra and T.~Schramm.
\href{https://doi.org/10.4230/LIPIcs.APPROX-RANDOM.2014.381}{Gap amplification for small-set expansion via random walks}.
APPROX/RANDOM 2014, LIPIcs 28, 381--391; arXiv:1310.1493.

\bibitem{rs10}
P.~Raghavendra and D.~Steurer.
\href{https://www.dsteurer.org/paper/expansion/}{Graph expansion and the unique games conjecture}.
STOC 2010, 755--764.

\bibitem{rst12}
P.~Raghavendra, D.~Steurer, and M.~Tulsiani.
\href{https://arxiv.org/abs/1011.2586}{Reductions between expansion problems}.
CCC 2012, 64--73; arXiv:1011.2586.

\bibitem{rvw02}
O.~Reingold, S.~Vadhan, and A.~Wigderson.
\href{https://arxiv.org/abs/math/0406038}{Entropy waves, the zig-zag graph product, and new constant-degree expanders}.
\emph{Ann. of Math.} 155 (2002), 157--187.

\bibitem{rs91}
N.~Robertson and P.~D.~Seymour.
Graph minors. X. Obstructions to tree-decomposition.
\emph{J. Combin. Theory Ser. B} 52 (1991), 153--190.

\bibitem{sv16}
S.~H.~S\ae{}ther and M.~Vatshelle.
\href{https://doi.org/10.1016/j.tcs.2015.11.039}{Hardness of computing width parameters based on branch decompositions over the vertex set}.
\emph{Theoret. Comput. Sci.} 615 (2016), 120--125.

\bibitem{sv95}
H.~Saran and V.~V.~Vazirani.
\href{https://ics.uci.edu/~vazirani/k-cut.pdf}{Finding \(k\) cuts within twice the optimal}.
\emph{SIAM J. Comput.} 24 (1995), 101--108.

\bibitem{sve11}
O.~Svensson.
\href{https://theory.epfl.ch/osven/Ola%20Svensson_publications/SICOMP11b.pdf}{Hardness of precedence constrained scheduling on identical machines}.
\emph{SIAM J. Comput.} 40 (2011), 1258--1274.

\bibitem{tz24}
A.~Ta-Shma and R.~Zadicario.
The expander hitting property when the sets are arbitrarily unbalanced.
APPROX/RANDOM 2024, LIPIcs 317, article 31,
\href{https://drops.dagstuhl.de/storage/00lipics/lipics-vol317-approx-random2024/LIPIcs.APPROX-RANDOM.2024.31/LIPIcs.APPROX-RANDOM.2024.31.pdf}{LIPIcs.APPROX-RANDOM.2024.31}; ECCC TR24-118.

\bibitem{tre13}
L.~Trevisan.
\href{https://arxiv.org/abs/1303.2730}{Is Cheeger-type approximation possible for nonuniform sparsest cut?}
arXiv:1303.2730.

\bibitem{vat12}
M.~Vatshelle.
New width parameters of graphs.
PhD thesis, University of Bergen, 2012.

\bibitem{wapl14}
Y.~Wu, P.~Austrin, T.~Pitassi, and D.~Liu.
\href{https://doi.org/10.1613/jair.4030}{Inapproximability of treewidth and related problems}.
\emph{J. Artificial Intelligence Res.} 49 (2014), 569--600.

\bibitem{wapl15}
Y.~Wu, P.~Austrin, T.~Pitassi, and D.~Liu.
\href{https://ijcai.org/proceedings/2015/600}{Inapproximability of treewidth and related problems (extended abstract)}.
IJCAI 2015, journal track.

\bibitem{yam18}
K.~Yamazaki.
\href{https://www.mdpi.com/1999-4893/11/11/173}{Inapproximability of rank, clique, Boolean, and maximum induced matching-widths under small set expansion hypothesis}.
\emph{Algorithms} 11 (2018), article 173.

\bibitem{zuc07}
D.~Zuckerman.
Linear degree extractors and the inapproximability of max clique and chromatic number.
\emph{Theory Comput.} 3 (2007), 103--128.

\end{thebibliography}

\end{document}
